\documentclass[10pt,a4paper]{article}
\usepackage[margin=2.2cm]{geometry}
\usepackage{fontspec}
\setmonofont{DejaVu Sans Mono}[Scale=0.84]
\newfontfamily\unifont{DejaVu Serif}[Scale=0.9]
\usepackage{amsmath,amssymb,mathtools}
\setcounter{MaxMatrixCols}{20}
\usepackage{longtable,booktabs,array}
\usepackage{fancyvrb}
\usepackage[dvipsnames]{xcolor}
\usepackage[colorlinks=true,linkcolor=NavyBlue,urlcolor=NavyBlue]{hyperref}
\usepackage{microtype}
\usepackage{needspace}
\usepackage{tocloft}\setlength{\cftbeforesecskip}{0.45em}
\setlength{\parindent}{0pt}\setlength{\parskip}{0.55em}
\setlength{\emergencystretch}{8em}
\makeatletter
\renewcommand\section{\@startsection{section}{1}{\z@}{-3.5ex \@plus -1ex \@minus -.2ex}{2.3ex \@plus.2ex}{\normalfont\Large\bfseries\raggedright}}
\renewcommand\subsection{\@startsection{subsection}{2}{\z@}{-3.25ex\@plus -1ex \@minus -.2ex}{1.5ex \@plus .2ex}{\normalfont\large\bfseries\raggedright}}
\renewcommand\subsubsection{\@startsection{subsubsection}{3}{\z@}{-3.25ex\@plus -1ex \@minus -.2ex}{1.5ex \@plus .2ex}{\normalfont\normalsize\bfseries\raggedright}}
\makeatother
\renewcommand{\arraystretch}{1.15}
\allowdisplaybreaks
\begin{document}
\title{The auxiliary torus of 7+1 spacetime: provenance\\[0.4em]\large The task, the decisions, the inputs, the data, the derivations and the verification record of the document on the auxiliary torus}
\author{Repository: github.com/once-ere/7plus1 — prepared for Nash\\ Notebook: 7+1.nb (Wolfram Language 15.0.1)\\ Document generated by tools/make\_docs.py on 2026-10-28}
\date{}
\maketitle
\begin{abstract}
This record states what was asked, how it was interpreted, every decision taken (among them the new section 25 of 7+1.nb, which computes every worked example), which results of the earlier sections and which data were used and where they come from, how each result was derived, and how it was checked: the 27 checks of section 25 (27 True), the 38 boolean checks of the independent Python program tools/ref\_torus\_brane.py (38 True), the comparison of the 17 groups of numbers printed by both programs, the check that the earlier sections of the notebook are unchanged, the program versions and the SHA-256 hashes of the files.
\end{abstract}
\tableofcontents
\Needspace{12\baselineskip}\section*{1. The task}\addcontentsline{toc}{section}{1. The task}
On 28 October 2026 the author asked for the following (summarized; the original wording is the author’s private working material and is not reproduced):

\begin{enumerate}
\item Define and explain the "auxiliary torus" precisely and correctly, in its own documentation file with its own symbol file and provenance file, each in Markdown and LaTeX with the LaTeX compiled to PDF.
\item Explain it for students; be clear, use examples, and organize the material well.
\item Do the same, in a separate set of three files, for the "world brane".
\item Push these files, and any other relevant files, to the repository github.com/once-ere/7plus1.
\item Then, in a second request: read a note of the author on the radion (a private PDF in the workspace, not part of the repository), use its ideas to verify and clarify the radion and brane documents in preparation, and create examples in which the four extra space dimensions are (1) a 4-sphere and (2) a 4-torus. For this document that means the comparison of the torus with a 4-sphere of the same volume (section 10 of the document) and the precise definition of the radion used in its sections 6.4–6.5 and 10.1.
\end{enumerate}
The standing rules of the project apply: every documentation file has a companion provenance file and a symbol file, in Markdown and LaTeX, compiled to PDF; no file refers the reader to another documentation, provenance or symbol file; every equation that is created or used is written in the file that uses it. The repository’s .gitignore is respected.

\textbf{Interpretation.} The auxiliary torus is the compact space of the four auxiliary coordinates $x_4,\dots,x_7$ of 7+1.nb, on which sections 17–24 rely (the Kaluza–Klein masses, the volume $V_\perp$, the radion, the Casimir energy). "Precisely" is read as: a definition as a quotient space, the metric quantities (size, volume) with their equations, its symmetries, and the consequences that the model uses, each stated as an equation; "correctly" as: every worked example is computed and checked twice, by the notebook and by an independent program. The coordinates of 7+1.nb are kept: $X=(x_0,\dots,x_7)$, $x_0$ the time (the direction of the octonion unit $e_0=1$), $x_1,x_2,x_3$ ordinary space, $x_4,\dots,x_7$ the four auxiliary space-like coordinates.

\Needspace{8\baselineskip}\subsection*{1.1 The signature: the author’s explanation of 27 October 2026}\addcontentsline{toc}{subsection}{1.1 The signature: the author’s explanation of 27 October 2026}
\textbf{Source.} On 27 October 2026 the author of the project explained that the xAct list $\{7,1,0\}$ declares seven space-like dimensions and one time-like dimension, while $\mathrm{diag}(+1,-1,\dots,-1)$ is simply the squares of the octonion basis, and asked that this be made clear in all documentation, examples and provenance records. \textbf{What it corrected.} The version of 7+1.nb pushed in the commits 52d8729 to dc2f178 of github.com/once-ere/7plus1 declared the xAct metric of section 12 with \texttt{DefMetric[\{1, 7, 0\}, …]}, reading the list as the signs of $\eta$; since then section 12 declares \texttt{DefMetric[\{7, 1, 0\}, g71[-a, -b], CD, …]}, exactly as the author’s notebook Einstein-Lovelock-7+1-Nash.nb does. \textbf{What was added.} Checks in 7+1.nb — section 4.7: $\eta$ is the squares of the octonion basis, $e_Ae_A=\eta_{AA}e_0$ and $\mathrm{Re}(e_Ae_B)=\eta_{AB}$; section 12.2: \texttt{SignatureOfMetric[g71]} $=\{7,1,0\}$ with \texttt{SignDetOfMetric[g71]} $=-1$, $\mathrm{Re}(e_Ae_B)=\eta_{AB}$, and $\mathrm{Re}(x\,x)=\eta_{AB}x^Ax^B$; section 17: the flat metric $\mathrm{diag}(+1,-1,-1,-1)$ of the world brane is the squares of the basis $e_0,e_1,e_2,e_3$ of the quaternion subalgebra — and the same statements in the Python programs \texttt{tools/\allowbreak{}ref71.py}, \texttt{tools/\allowbreak{}ref\_\allowbreak{}gravity.py} and \texttt{tools/\allowbreak{}ref\_\allowbreak{}sm.py}; the explanation itself in the text of 7+1.nb (sections 1.2, 1.4, 4.2, 4.7, 7, 11, 12.2, 13, 15.1, 17) and in every document, symbol file and provenance file of the repository. \textbf{What did not change.} No result. Compared output by output with the backup below, every output of 7+1.nb is the same, except the declared list itself, the new checks and the check counts (and, as in every run, the name of the private context and the measured times); in particular every component of every Lovelock tensor is unchanged. The components follow from the component matrix of the metric, and the sign of the determinant, \texttt{SignDetOfMetric}, is $-1$ for both lists. \textbf{Backup.} The evaluated notebook as it was before is kept in the workspace as \texttt{backups/\allowbreak{}7+1\_\allowbreak{}2026-10-27\_\allowbreak{}before\_\allowbreak{}signature.nb} (the folder backups/ is excluded from the repository by .gitignore), 1034401 bytes, SHA-256 \texttt{06094248\allowbreak{}8caf4f45\allowbreak{}120ffd37\allowbreak{}be0123e9\allowbreak{}6992a31c\allowbreak{}d7fdbde0\allowbreak{}9ee4dc4c\allowbreak{}ed4f5d3a}; \texttt{tools/\allowbreak{}verify\_\allowbreak{}7p1\_\allowbreak{}nb.py} compares every input, text and section cell of 7+1.nb with it and finds every cell present, each one unchanged or changed by this explanation, and it compares every output in the same way.

\textbf{The explanation, in full.} The 7+1 spacetime has one time-like dimension ($x_0$, the direction of the octonion unit $e_0=1$) and seven space-like dimensions ($x_1,x_2,x_3$ ordinary space, $x_4,\dots,x_7$ auxiliary space). Two different objects describe this, and neither may be replaced by the other. (i) The tensor package xAct, used in section 12 of 7+1.nb exactly as in the author’s notebook Einstein-Lovelock-7+1-Nash.nb, declares the metric with \texttt{DefMetric[\{7, 1, 0\}, g71[-a, -b], CD, …]}: the list $\{7,1,0\}$ declares seven space-like dimensions, one time-like dimension and no degenerate dimension — the numbers of entries $+1$ and $-1$ of $\eta_{71}=\mathrm{diag}(-1,+1,\dots,+1)$ (\texttt{η71}), the same spacetime written in the sign convention of 4+4.nb (space-like $=+1$, time-like $=-1$). (ii) The matrix $\eta=-\eta_{71}=\mathrm{diag}(+1,-1,\dots,-1)$ (\texttt{η17}) is simply the squares of the octonion basis: $e_0^2=+1$ for the unit $e_0=1$ (time) and $e_i^2=-1$ for the seven imaginary units $e_1,\dots,e_7$ (space); more generally it is the real part of the product of two basis units, and it gives the quadratic form of the octonions:
\begin{gather*}
\eta_{AB}=\mathrm{Re}(e_Ae_B),\qquad e_0^2=+1,\qquad e_i^2=-1\ \ (i=1,\dots,7),\\ \mathrm{Re}(x\,x)=x_0^2-x_1^2-\cdots-x_7^2=\eta_{AB}\,x^Ax^B\qquad (x=x^Ae_A\in\mathbb O).
\end{gather*}
$\eta$ gives the signs of the Clifford relation $\Gamma_A\Gamma_B+\Gamma_B\Gamma_A=2\,\eta_{AB}\,1_{16}$ of the real 16×16 Dirac matrices: time squares to $+1_{16}$ and each of the seven space directions to $-1_{16}$, exactly like the octonion units (in the notation $Cl(p,q)$, with $p$ generators squaring to $+1$ and $q$ to $-1$, the algebra $Cl(1,7)\cong\mathbb R(16)$). It also gives the signs of the metric components: $g_{00}>0$ along the time-like direction, negative squares along the space-like directions. Every component computed by 7+1.nb follows from the component matrix of the metric; the declared list fixes for xAct the sign of the determinant of the metric, \texttt{SignDetOfMetric[g71]} $=-1$, which is the same in both sign conventions, since $\det \eta=\det \eta_{71}=-1$. 7+1.nb checks \texttt{SignatureOfMetric[g71]} $=\{7,1,0\}$, \texttt{SignDetOfMetric[g71]} $=-1$, $\mathrm{Re}(e_Ae_B)=\eta_{AB}$ and $\mathrm{Re}(x\,x)=\eta_{AB}x^Ax^B$.

\Needspace{12\baselineskip}\section*{2. Decisions}\addcontentsline{toc}{section}{2. Decisions}
{\footnotesize
\begin{longtable}{>{\raggedright\arraybackslash}p{\dimexpr 0.0500\linewidth-2\tabcolsep\relax}>{\raggedright\arraybackslash}p{\dimexpr 0.4700\linewidth-2\tabcolsep\relax}>{\raggedright\arraybackslash}p{\dimexpr 0.4800\linewidth-2\tabcolsep\relax}}
\toprule
\textbf{\#} & \textbf{decision} & \textbf{reason / equation} \\
\midrule\endfirsthead
\toprule
\textbf{\#} & \textbf{decision} & \textbf{reason / equation} \\
\midrule\endhead
\endfoot
\bottomrule\endlastfoot
1 & append section 25 to 7+1.nb ("The auxiliary torus and the world brane: definitions and worked examples"); the summary of all checks becomes section 26 & every worked example of the document is then computed and checked by the notebook itself, in the same evaluation as everything else \\
2 & write an independent Python program, tools/ref\_torus\_brane.py, for the same examples & its own lattice sieve, Fourier transform (mpmath), symbolic integrals and curvature (sympy); no code shared with the notebook \\
3 & back up the evaluated notebook before appending section 25 and check that sections 0–24 are unchanged & tools/verify\_7p1\_nb.py compares every cell and output with the backup (section 6.3) \\
4 & the torus of 7+1.nb is the cubic torus $\mathbb R^4/(L_c\mathbb Z)^4$ with periodic fermions & the choice of sections 18 and 22–24; other lattices and spin structures are named but not used \\
5 & keep both $L_c$ and $b$, and state that only $bL_c$ is physical & 7+1.nb uses both symbols; $\vec y=\lambda\vec y\,'$ maps $(L_c,b)\to(L_c/\lambda,\lambda b)$ \\
6 & the size is defined by $b^8=\det(-g_{kl})$ & the definition of 7+1.nb (section 20.1); for the metrics of the model $-g_{kl}=b^2\delta_{kl}$ \\
7 & the Newton-law example in linearized gravity without the stabilizing potential; the stabilizer is added through the radion mass & the simplest setting in which images and Kaluza–Klein modes can be compared exactly \\
8 & the image sum is summed over the shells $k\le k_{\max}$ of section 22.2 and closed by a continuum tail; the Kaluza–Klein sum over the shells $k\le20000$ & tail and truncation errors are far below the tolerance $10^{-8}$ (section 4) \\
9 & the Newtonian potential is written $\Phi_{\rm N}$ (not $\Phi$), the PDG radius $R_{\rm PDG}$ (not $R$) & $\Phi$ is the Higgs doublet and $R$ the scalar curvature in 7+1.nb \\
10 & numbers today with $M_8=10$ TeV, $M_{\rm Pl}=2.435\times10^{18}$ GeV, $m_r=10$ GeV & the choices of sections 22–24 \\
11 & the check label "7.7 10\textasciicircum{}17 states below m\_h" of section 24.4 is left unchanged; the document states its exact meaning & it is the weighted sum $\sum(1-\mu_{\vec n}^2/m_h^2)^2=V_\perp m_h^4/(192\pi^2)$, not the number $N(m_h)=V_\perp m_h^4/(32\pi^2)$ of levels; changing a label of section 24 would change a verified section \\
12 & the orbits of the shells are computed with $|n_i|\le3$ & enough for $k\le12$, since $|n_i|^2\le k$ \\
13 & compare the torus with a 4-sphere of the SAME volume $V_\perp$ (section 26 of 7+1.nb, appended for the second request) & the same four-dimensional Newton constant $8\pi G_4=\kappa/V_\perp$; the differences are then those of the geometry alone \\
14 & the radion in the Einstein frame, $ds^2=e^{-4\beta}g_E-e^{2\beta}\gamma$, on both spaces & the conventions of section 20.3 of 7+1.nb; the reduction is checked in section 26.1 \\*
\end{longtable}}
\Needspace{12\baselineskip}\section*{3. Inputs from the earlier sections of 7+1.nb}\addcontentsline{toc}{section}{3. Inputs from the earlier sections of 7+1.nb}
{\footnotesize
\begin{longtable}{>{\raggedright\arraybackslash}p{\dimexpr 0.2400\linewidth-2\tabcolsep\relax}>{\raggedright\arraybackslash}p{\dimexpr 0.2400\linewidth-2\tabcolsep\relax}>{\raggedright\arraybackslash}p{\dimexpr 0.1000\linewidth-2\tabcolsep\relax}>{\raggedright\arraybackslash}p{\dimexpr 0.4200\linewidth-2\tabcolsep\relax}}
\toprule
\textbf{object} & \textbf{Mathematica} & \textbf{section} & \textbf{equation or value} \\
\midrule\endfirsthead
\toprule
\textbf{object} & \textbf{Mathematica} & \textbf{section} & \textbf{equation or value} \\
\midrule\endhead
\endfoot
\bottomrule\endlastfoot
cosmological metric & \texttt{gCos}, \texttt{a3[x0]}, \texttt{b4[x0]}, \texttt{cvCos} & 11.2 & $ds^2=dx_0^2-a^2d\vec x^2-b^2d\vec y^2$; $R=-2(3H_a'+6H_a^2+4H_b'+10H_b^2+12H_aH_b)$ \\
inhomogeneous metric & \texttt{gPoly}, \texttt{cvPoly} & 11.3 & $n=1+\frac{x_0x_4}5+\frac{x_0^2}{11}$, $a=1+\frac{x_0}3+\frac{x_4^2}7$, $b=2-\frac{x_0x_4}4+\frac{x_0^2}9+\frac{x_4}{13}$ \\
Dirac matrices, brane representation, intertwiner & \texttt{Γ17}, \texttt{g4D}, \texttt{g5D}, \texttt{rhoK}, \texttt{rBr}, \texttt{uInt} & 3–5, 17.2 & $\Gamma_AU=UR_A$, $R_\mu=\gamma_\mu\otimes1_4$, $R_k=\gamma_5\otimes\rho_k$ \\
Kaluza–Klein operator & \texttt{xKK} & 18.1 & $X=\gamma_5\otimes(\vec\rho\cdot\vec q)/b_0$, $X^2=-|\vec q|^2/b_0^2$ \\
stress tensor from a comoving energy & \texttt{feStress[W]} & 20.2 & $\varepsilon=W/(a^3b^4)$, $p_a=-\partial_aW/(3a^2b^4)$, $p_b=-\partial_bW/(4a^3b^3)$ \\
radion, stabilizing potential, four-dimensional constants & \texttt{uPot}, \texttt{kU} & 20.3 & $U=U_0+\frac{K_U}2\beta^2$, $m_r^2=\kappa K_U/12$, $8\pi G_4=\kappa/V_\perp$, $\phi=\sqrt{12}M_{\rm Pl}\beta$ \\
Kasner exponents & \texttt{rIp1}, \texttt{rIq1}, \texttt{rIp2}, \texttt{rIq2} & 21.1 & $p_1=(1+2\sqrt2)/7$, $q_1=(2-3\sqrt2)/14$, $3p_2+4q_2=3$ \\
cubic regime & \texttt{rIcub}, \texttt{rIhStar} & 21.2–21.3 & $h_0=-\sqrt{\alpha_2/(30\alpha_3)}$, $h_*=-\sqrt{\alpha_2/(54\alpha_3)}$ \\
decompactification without stabilizer & \texttt{rIlate} & 21.3 & $\tau H_b\to\frac25$, $p_b/\varepsilon\to\frac14$ \\
capture of the radion & \texttt{rIIaF}, \texttt{rIIphF}, \texttt{rIItEF} & 22.1 & $\beta\simeq A\sin(m_r\tau+\vartheta)/(m_r(\tau-\tau_E))$ \\
Jacobi’s $r_4(k)$, the volume today & \texttt{rIIr4}, \texttt{rIIsig}, \texttt{rIIkmax}, \texttt{rIIVp} & 22.2 & $r_4=8\sigma(k)-32\sigma(k/4)$; $V_\perp=M_{\rm Pl}^2/M_8^6$ \\
the radion today & — & 22.4 & $\beta_{\rm eq}$ from $\Lambda_4$ and matter, from a tension $f_W=M_8^4$ \\
the far future & \texttt{rIIIH} & 23.1–23.2 & $H_\infty=H_0\Omega_\Lambda^{1/2}$; eight-dimensional de Sitter $21H^2-420\alpha_2H^4+2520\alpha_3H^6=\Lambda$ \\
tower sums, Casimir energy, $\rho_\Lambda$ & \texttt{dkTowN}, \texttt{dkCas}, \texttt{dkRhoL} & 24.4–24.5 & $\sum(1-\mu^2/m_h^2)^2=V_\perp m_h^4/(192\pi^2)$; $\rho_{\rm Cas}$; $3H_0^2M_{\rm Pl}^2\Omega_\Lambda$ \\
the radion on $T^4$ and $S^4$, the sphere today & \texttt{rdV}, \texttt{rdDl}, \texttt{rdZon}, \texttt{rdKS}, \texttt{rdRS4}, \texttt{rdToday} & 26.1–26.4 & $V(\beta)=M_{\rm Pl}^2[e^{-4\beta}(\Lambda+\kappa U)-6\sigma e^{-6\beta}/r_0^2]$; $d_l=(2l+3)(l+1)(l+2)/6$; $r_{S^4}=(3V_\perp/(8\pi^2))^{1/4}$ \\*
\end{longtable}}
Before section 25 was appended, the evaluated notebook (sections 0–24 and the summary of all checks, then numbered 25) was copied to \texttt{backups/\allowbreak{}7+1\_\allowbreak{}2026-10-28\_\allowbreak{}before\_\allowbreak{}torus\_\allowbreak{}brane.nb} (1940664 bytes, SHA-256 \texttt{39f64ee8\allowbreak{}1d35c0c9\allowbreak{}fb94188f\allowbreak{}d7f8c8e7\allowbreak{}5c96ea19\allowbreak{}8fbf72a5\allowbreak{}ca6b3c00\allowbreak{}27426018}; excluded from the repository by .gitignore).

\Needspace{12\baselineskip}\section*{4. Data and their sources}\addcontentsline{toc}{section}{4. Data and their sources}
{\footnotesize
\begin{longtable}{>{\raggedright\arraybackslash}p{\dimexpr 0.2200\linewidth-2\tabcolsep\relax}>{\raggedright\arraybackslash}p{\dimexpr 0.3400\linewidth-2\tabcolsep\relax}>{\raggedright\arraybackslash}p{\dimexpr 0.4400\linewidth-2\tabcolsep\relax}}
\toprule
\textbf{quantity} & \textbf{value used} & \textbf{source} \\
\midrule\endfirsthead
\toprule
\textbf{quantity} & \textbf{value used} & \textbf{source} \\
\midrule\endhead
\endfoot
\bottomrule\endlastfoot
$\hbar c$ & $1.973269804\times10^{-16}$ GeV m & CODATA 2018 (E. Tiesinga et al., Rev. Mod. Phys. 93 (2021) 025010) \\
proton charge radius $r_p$ & 0.8414 fm & CODATA 2018 \\
Bohr radius $a_0$ & $5.29177210903\times10^{-11}$ m & CODATA 2018 \\
$M_{\rm Pl}$ (reduced) & $2.435\times10^{18}$ GeV & $(8\pi G_N)^{-1/2}$, as in sections 17–24 \\
$M_8$, $m_r$ & 10 TeV, 10 GeV & the choices of sections 22–24 (natural values) \\
$m_h$ & 125.20 GeV & Particle Data Group 2025 (as in section 17.3) \\
bounds on $M_D$ & 5.9 to 11.2 TeV for 6 to 2 extra dimensions & Particle Data Group, Phys. Rev. D 110 (2024) 030001 (ATLAS, 139 fb$^{-1}$) \\
$\rho_\Lambda$ & $2.518\times10^{-47}$ GeV$^4$ & $3H_0^2M_{\rm Pl}^2\Omega_\Lambda$ with the Planck 2018 values of section 22.4 \\
$\alpha_3/\alpha_2^2$ & $10^{-6}$ & the choice of sections 21–22.1 \\*
\end{longtable}}
Software: Wolfram Language 15.0.1 (front end WolframNB.exe; the kernel configuration "Eval71" is a fresh kernel); Python 3.13.17 with numpy 2.5.3, scipy 1.18.1, sympy 1.14.0, mpmath 1.3.0.

\Needspace{12\baselineskip}\section*{5. Derivations}\addcontentsline{toc}{section}{5. Derivations}
\textbf{The volume.} For $-g_{kl}=b^2\delta_{kl}$, $\det(-g_{kl})=b^8$ and $V_\perp=\int_{[0,L_c)^4}\sqrt{\det(-g_{kl})}\,d^4y=b^4L_c^4$; under $\vec y=\lambda\vec y\,'$ the line element $b^2d\vec y\cdot d\vec y=(\lambda b)^2d\vec y\,'\cdot d\vec y\,'$, so $(L_c,b)\to(L_c/\lambda,\lambda b)$ and $bL_c$, $b^4L_c^4$ are unchanged.

\textbf{The orbits.} A signed permutation permutes the entries of $\vec n$ and changes their signs; the orbit of $\vec n$ is the set of vectors with the same multiset $\{|n_i|\}$. Its size is the number of distinct arrangements of the multiset times $2^{z}$, $z$ the number of non-zero entries:
\begin{equation*}
|\text{orbit}(\vec n)|=\frac{4!}{\prod_jm_j!}\,2^{z},\qquad\text{e.g. }(1,1,0,0):\ \frac{4!}{2!\,2!}\cdot2^2=24,\quad (1,1,1,1):\ 1\cdot2^4=16 .
\end{equation*}
\textbf{Orthonormality and Parseval.} $\int_0^{L_c}e^{2\pi i(m-n)y/L_c}dy=L_c\delta_{mn}$, so $\int_Fb^4\phi^*_{\vec n}\phi_{\vec m}=b^4L_c^4\delta_{\vec n\vec m}/V_\perp=\delta_{\vec n\vec m}$. With $\cos^2\theta=\frac12+\frac14e^{2i\theta}+\frac14e^{-2i\theta}$ the bump $f=\prod\cos^2(\pi y_k/L_c)$ has $c_{\vec n}=\sqrt{V_\perp}\prod w(n_k)$, and
\begin{equation*}
\sum_{\vec n}|c_{\vec n}|^2=V_\perp\Big(\frac14+\frac1{16}+\frac1{16}\Big)^4=\Big(\frac38\Big)^4V_\perp,\qquad \int_Fb^4f^2\,d^4y=b^4\Big(\int_0^{L_c}\cos^4\frac{\pi y}{L_c}dy\Big)^4=b^4\Big(\frac38L_c\Big)^4 .
\end{equation*}
\textbf{Kaluza–Klein masses.} For $g=\mathrm{diag}(1,-1,-1,-1,-b_0^2,\dots,-b_0^2)$, $\Box=\partial_0^2-\nabla_{\vec x}^2-b_0^{-2}\nabla_{\vec y}^2$; on $e^{-iEx_0+i\vec k\cdot\vec x+i\vec q\cdot\vec y}$ this is $-E^2+|\vec k|^2+|\vec q|^2/b_0^2$, so $E^2=|\vec k|^2+m^2+|\vec q|^2/b_0^2$ and $\mu_{\vec n}=|\vec q_{\vec n}|/b_0=2\pi|\vec n|/(b_0L_c)=2\pi|\vec n|/\ell_\perp$. For dirac16complex, $(\gamma\cdot p\otimes1_4+X)^2=(\gamma\cdot p)^2\otimes1_4+X^2$ because $\gamma_5$ anticommutes with every $\gamma_\mu$, and $X^2=\gamma_5^2\otimes(\vec\rho\cdot\vec q)^2/b_0^2=-|\vec q|^2/b_0^2$ from $\rho_k\rho_l+\rho_l\rho_k=-2\delta_{kl}$.

\textbf{The graviton states.} A massless particle in $D$ dimensions has the little group $SO(D-2)$; the graviton is its symmetric traceless tensor, $\frac{(D-2)(D-1)}2-1=\frac{D(D-3)}2=20$ states. At the zero mode the $SO(2)$ of four dimensions sees $g_{\mu\nu}$ (2), four vectors $g_{\mu k}$ ($4\times2$) and the symmetric $g_{kl}$ (10). At a massive level with $\vec q\ne0$ the little group of the four-dimensional massive particle is $SO(3)$, and the 20 states form a massive spin 2 (5), three massive vectors (the torus directions orthogonal to $\vec q$; $3\times3$) and the 6 components of a symmetric $3\times3$ tensor.

\textbf{Jacobi and the lattice sums.} $\sum_{k\ge1}\sigma(k)k^{-s}=\zeta(s)\zeta(s-1)$ and $\sum_k\sigma(k/4)k^{-s}=4^{-s}\zeta(s)\zeta(s-1)$, so
\begin{equation*}
\sum_{k\ge1}\frac{r_4(k)}{k^s}=8(1-4^{1-s})\,\zeta(s)\,\zeta(s-1):\qquad Z(\tfrac52)=7\,\zeta(\tfrac52)\,\zeta(\tfrac32),\qquad Z(4)=8(1-4^{-3})\,\zeta(4)\,\zeta(3).
\end{equation*}
\textbf{Weyl’s law and the level count.} On average $r_4(k)\approx\frac{d}{dk}\big(\frac{\pi^2}2k^2\big)=\pi^2k$, so $N(R)\simeq\frac{\pi^2}2R^4$; a level below $\mu$ has $|\vec n|\le R=\mu\ell_\perp/(2\pi)$:
\begin{equation*}
N(\mu)\simeq\frac{\pi^2}2\Big(\frac{\mu\ell_\perp}{2\pi}\Big)^4=\frac{V_\perp\mu^4}{32\pi^2},\qquad \sum_{\mu_{\vec n}<m_h}\Big(1-\frac{\mu_{\vec n}^2}{m_h^2}\Big)^2\simeq\frac{V_\perp m_h^4}{8\pi^2}\int_0^1u^3(1-u^2)^2du=\frac{V_\perp m_h^4}{192\pi^2}.
\end{equation*}
\textbf{From eight to four dimensions.} For $g=\mathrm{diag}(g^{(4)}_{\mu\nu}(x),-b_0^2\delta_{kl})$ the Christoffel symbols with an index $k$ vanish (nothing depends on $\vec y$ and $b_0$ is constant), so $R=R^{(4)}$ and $\sqrt{|g|}=b_0^4\sqrt{|g^{(4)}|}$; integrating $-\frac1{2\kappa}\int\sqrt{|g|}R$ over the torus gives $-\frac{V_\perp}{2\kappa}\int d^4x\sqrt{|g^{(4)}|}R^{(4)}=-\frac1{16\pi G_4}\int d^4x\sqrt{|g^{(4)}|}R^{(4)}$, i.e. $8\pi G_4=\kappa/V_\perp$.

\textbf{Newton’s law in $D$ dimensions.} The linearized Einstein equation for a static source gives $\nabla^2\Phi_{\rm N}=8\pi G_D\frac{D-3}{D-2}\rho$; with $\nabla^2r^{-(d-2)}=-(d-2)\Omega_{d-1}\delta^d$ ($d=D-1$, $\Omega_{d-1}=2\pi^{d/2}/\Gamma(d/2)$):
\begin{align*}
\Phi_{\rm N}&=-\frac{8\pi G_D(D-3)M}{(D-2)(D-3)\Omega_{D-2}r^{D-3}}=-\frac{\kappa M}{(D-2)\Omega_{D-2}r^{D-3}}: \\
D=8:&\ \ -\frac{\kappa M}{6\cdot\frac{16\pi^3}{15}r^5}=-\frac{5\kappa M}{32\pi^3r^5},\qquad D=4:\ \ -\frac{GM}r .
\end{align*}
\textbf{Images and Kaluza–Klein modes.} The Poisson summation formula $\sum_{\vec n}f(\vec n)=\sum_{\vec m}\hat f(2\pi\vec m)$ with $\hat f(\vec p)=\int d^4y\,e^{-i\vec p\cdot\vec y}f(\vec y)$, and for a radial function in four dimensions
\begin{equation*}
\hat f(p)=\frac{4\pi^2}p\int_0^\infty y^2J_1(py)f(y)\,dy .
\end{equation*}
For $f=(r^2+y^2)^{-5/2}$ the integral is the case $\nu=1$, $\mu=\frac32$ of
\begin{equation*}
\int_0^\infty\frac{y^{\nu+1}J_\nu(py)}{(y^2+r^2)^{\mu+1}}dy=\frac{p^\mu r^{\nu-\mu}}{2^\mu\Gamma(\mu+1)}K_{\nu-\mu}(pr)\quad(-1<\nu<2\mu+\tfrac32),
\end{equation*}
with $\Gamma(\frac52)=\frac34\sqrt\pi$ and $K_{-1/2}(z)=K_{1/2}(z)=\sqrt{\pi/(2z)}\,e^{-z}$:
\begin{equation*}
\frac{p^{3/2}r^{-1/2}}{2^{3/2}\cdot\frac34\sqrt\pi}\sqrt{\frac\pi{2pr}}\,e^{-pr}=\frac p{3r}e^{-pr},\qquad \hat f(p)=\frac{4\pi^2}p\cdot\frac p{3r}e^{-pr}=\frac{4\pi^2}{3r}e^{-pr},
\end{equation*}
hence $\sum_{\vec n}(r^2+|\vec n|^2)^{-5/2}=\frac{4\pi^2}{3r}\sum_{\vec m}e^{-2\pi r|\vec m|}$. \textbf{The numerical evaluation.} The left side is summed over the shells $k\le K=k_{\max}$ and closed with the continuum tail ($r_4(k)\to\pi^2k$, $u_0=r^2+K+\frac12$)
\begin{equation*}
\sum_{k>K}\frac{r_4(k)}{(r^2+k)^{5/2}}\simeq\pi^2\int_{u_0}^\infty\frac{t-r^2}{t^{5/2}}dt=\pi^2\Big(\frac2{\sqrt{u_0}}-\frac{2r^2}{3u_0^{3/2}}\Big),\qquad \sum_{k>K}\frac{r_4(k)}{k^{5/2}}\simeq\frac{2\pi^2}{\sqrt{K+\frac12}},
\end{equation*}
and the right side over $k\le20000$ (the omitted terms are below $e^{-2\pi r\sqrt{20000}}<10^{-96}$ for $r\ge\frac14$). Near the source $S(r)=r^{-5}+\sum_{\vec n\ne0}(r^2+|\vec n|^2)^{-5/2}=r^{-5}+Z(\frac52)+O(r^2)$; far from it $\frac{3r}{4\pi^2}S(r)=1+8e^{-2\pi r}+24e^{-2\sqrt2\pi r}+\dots$

\textbf{The factor $\frac53$.} Restoring $\ell_\perp$: $\Phi_{\rm N}\to-\frac{5\kappa M}{32\pi^3}\cdot\frac{4\pi^2}{3V_\perp r}=-\frac{5\cdot8\pi G_4M}{24\pi r}=-\frac53\frac{G_4M}r$. In terms of exchanged particles: one-graviton exchange between two sources in $D$ dimensions is proportional to $T_{AB}T'^{AB}-T\,T'/(D-2)$; for two masses at rest ($T_{00}=M$, $T=M$) this is $MM'(D-3)/(D-2)$: $\frac56MM'$ for the zero mode of eight-dimensional gravity, $\frac12MM'$ for the four-dimensional graviton. The ratio is $\frac53$; the difference, $\frac13MM'=\frac23\cdot\frac12MM'$, is the radion, whose coupling $\phi T^\mu{}_\mu/(\sqrt3M_{\rm Pl})$ gives the potential $-\frac{MM'}{3M_{\rm Pl}^2}\frac1{4\pi r}=-\frac23\frac{G_4MM'}r$ ($8\pi G_4=M_{\rm Pl}^{-2}$).

\textbf{The radion coupling.} With $g_E=(b/b_0)^4h$ the four-dimensional Einstein frame (the Planck mass is $\propto b^4$), $h=e^{-4\beta}g_E$; varying $h^{\mu\nu}=e^{4\beta}g_E^{\mu\nu}$ by $4\beta g_E^{\mu\nu}$ in $\delta S=\frac12\int\sqrt{|h|}T_{\mu\nu}\delta h^{\mu\nu}$ gives $2\beta T^\mu{}_\mu$, and $\beta=\phi/(\sqrt{12}M_{\rm Pl})$ gives $\phi T^\mu{}_\mu/(\sqrt3M_{\rm Pl})$ (section 20.3 of 7+1.nb).

\textbf{The radion today.} Static and linear in $\beta$, the $44$ equation of section 20.3, $\frac\kappa4bU'(b)=\Lambda+\kappa U+2\kappa f_W/V_\perp+\frac\kappa2(\varepsilon-3p_a)+\kappa p_b-12\alpha_2H_a^2(H_a'+H_a^2)$, with $bU'=K_U\beta$ and $\frac{\kappa K_U}4=3m_r^2$, gives $\beta_{\rm eq}$ of the document (with $\Lambda_4=\Lambda+\kappa U_0+\kappa f_W/V_\perp$).

\textbf{The delta function and smearing.} On each circle the integral $\int_{-L_c/2}^{L_c/2}e^{-2\pi iny/L_c}\delta(y)dy=1$ gives the Fourier series $\delta(y)=L_c^{-1}\sum_ne^{2\pi iny/L_c}$; the zero mode of $\delta^4(\vec y)$ is $L_c^{-4}$, and $\frac1{V_\perp}\int_Fb^4\frac{\rho_4\delta^4}{b^4}=\frac{\rho_4}{V_\perp}$.

\Needspace{12\baselineskip}\section*{6. Verification}\addcontentsline{toc}{section}{6. Verification}
\Needspace{8\baselineskip}\subsection*{6.1 The checks of section 25 of 7+1.nb}\addcontentsline{toc}{subsection}{6.1 The checks of section 25 of 7+1.nb}
\textbf{Evaluation of 7+1.nb}: Wolfram 15.0.1 for Microsoft Windows (64-bit) (July 2, 2026), a fresh kernel (front-end kernel configuration Eval71); finished Wed 28 Oct 2026 08:57:47; 283.883 s; message cells 0; output cells 554. The checks of section 25 (checks 1–14 concern the torus alone, 15–16 the torus and the brane together, 17–27 the world brane):

{\footnotesize
\begin{longtable}{>{\raggedright\arraybackslash}p{\dimexpr 0.0550\linewidth-2\tabcolsep\relax}>{\raggedright\arraybackslash}p{\dimexpr 0.7450\linewidth-2\tabcolsep\relax}>{\raggedright\arraybackslash}p{\dimexpr 0.0800\linewidth-2\tabcolsep\relax}>{\raggedright\arraybackslash}p{\dimexpr 0.1200\linewidth-2\tabcolsep\relax}}
\toprule
\textbf{\#} & \textbf{check} & \textbf{result} & \textbf{section} \\
\midrule\endfirsthead
\toprule
\textbf{\#} & \textbf{check} & \textbf{result} & \textbf{section} \\
\midrule\endhead
\endfoot
\bottomrule\endlastfoot
1 & \texttt{auxiliary torus: for the cosmological metric det(-g\_\allowbreak{}kl) == b\textasciicircum{}8 (k, l = 4..7), and the proper volume of the fundamental domain [0, Lc)\textasciicircum{}4, Integrate[Sqrt[det(-g\_\allowbreak{}kl)], \{y, 0, Lc\}\textasciicircum{}4], is b\textasciicircum{}4 Lc\textasciicircum{}4 = V\_\allowbreak{}perp} & True & 25.1 \\
2 & \texttt{only the product b Lc of the auxiliary torus is physical: y = lambda y' maps the torus of period Lc and size b to the one of period Lc/\allowbreak{}lambda and size lambda b; the line element b\textasciicircum{}2 dy.dy, the proper period b Lc and V\_\allowbreak{}perp = b\textasciicircum{}4 Lc\textasciicircum{}4 are unchanged} & True & 25.1 \\
3 & \texttt{symmetries of the auxiliary torus that fix y = 0 (and the brane): the 384 = 2\textasciicircum{}4 4! signed permutation matrices are distinct orthogonal integer matrices (they map Z\textasciicircum{}4 onto itself and preserve |n|\textasciicircum{}2); the shells |n|\textasciicircum{}2 = 1, 2, 3, 4, 9 split into orbits of sizes \{8\}, \{24\}, \{32\}, \{8, 16\}, \{8, 96\}} & True & 25.1 \\
4 & \texttt{Fourier modes of the auxiliary torus: phi\_\allowbreak{}n = Exp[I q\_\allowbreak{}n.y]/\allowbreak{}Sqrt[V\_\allowbreak{}perp] with q\_\allowbreak{}n = 2 Pi n/\allowbreak{}Lc (n in Z\textasciicircum{}4) are orthonormal for the proper volume element b\textasciicircum{}4 d\textasciicircum{}4y (exact integrals: n = m = (1, 0, -2, 3); n = (1, 0, -2, 3), m = (1, 1, -2, 3); n = 0, m = (0, 2, 0, 0))} & True & 25.2 \\
5 & \texttt{Parseval on the auxiliary torus: f(y) = Prod\_\allowbreak{}k cos\textasciicircum{}2(Pi y\_\allowbreak{}k/\allowbreak{}Lc) = Prod\_\allowbreak{}k (1/\allowbreak{}2 + (1/\allowbreak{}4) e\textasciicircum{}(2 Pi I y\_\allowbreak{}k/\allowbreak{}Lc) + (1/\allowbreak{}4) e\textasciicircum{}(-2 Pi I y\_\allowbreak{}k/\allowbreak{}Lc)) has the 81 non-zero coefficients c\_\allowbreak{}n = Sqrt[V\_\allowbreak{}perp] Prod\_\allowbreak{}k w(n\_\allowbreak{}k), w(0) = 1/\allowbreak{}2, w(+-1) = 1/\allowbreak{}4 (three of them checked by integration), and Sum\_\allowbreak{}n |c\_\allowbreak{}n|\textasciicircum{}2 == Int b\textasciicircum{}4 d\textasciicircum{}4y f\textasciicircum{}2 == (3/\allowbreak{}8)\textasciicircum{}4 V\_\allowbreak{}perp} & True & 25.2 \\
6 & \texttt{the Laplacian of the auxiliary torus: -(1/\allowbreak{}b\textasciicircum{}2) Sum\_\allowbreak{}k d\textasciicircum{}2 phi\_\allowbreak{}n/\allowbreak{}dy\_\allowbreak{}k\textasciicircum{}2 == (|q\_\allowbreak{}n|\textasciicircum{}2/\allowbreak{}b\textasciicircum{}2) phi\_\allowbreak{}n (n = (2, -1, 0, 3))} & True & 25.2 \\
7 & \texttt{scalar Kaluza-Klein masses: for the static metric diag(1, -1, -1, -1, -b0\textasciicircum{}2, ..., -b0\textasciicircum{}2) the Klein-Gordon operator Box + m\textasciicircum{}2 (Box f = |g|\textasciicircum{}(-1/\allowbreak{}2) d\_\allowbreak{}A(|g|\textasciicircum{}(1/\allowbreak{}2) g\textasciicircum{}AB d\_\allowbreak{}B f)) on Exp[-I E x0 + I k.x + I q.y] gives (-E\textasciicircum{}2 + |k|\textasciicircum{}2 + |q|\textasciicircum{}2/\allowbreak{}b0\textasciicircum{}2 + m\textasciicircum{}2) times the wave: in ordinary spacetime a particle of mass Sqrt[m\textasciicircum{}2 + |q|\textasciicircum{}2/\allowbreak{}b0\textasciicircum{}2]} & True & 25.2 \\
8 & \texttt{Jacobi's four-square theorem: the number of n in Z\textasciicircum{}4 with |n|\textasciicircum{}2 = k is r\_\allowbreak{}4(k) = 8 sigma(k) - 32 sigma(k/\allowbreak{}4) for k = 1..100 (direct count over \{-10, ..., 10\}\textasciicircum{}4)} & True & 25.3 \\
9 & \texttt{Weyl's law: the number N(R) of n in Z\textasciicircum{}4 with |n| <= R approaches the volume Pi\textasciicircum{}2 R\textasciicircum{}4/\allowbreak{}2 of the four-ball: N(10) = 49689, N(100) = 493490641, and |N(R)/\allowbreak{}(Pi\textasciicircum{}2 R\textasciicircum{}4/\allowbreak{}2) - 1| < 10\textasciicircum{}-2, 10\textasciicircum{}-3, 3 10\textasciicircum{}-5, 5 10\textasciicircum{}-6, 10\textasciicircum{}-6 for R = 10, 30, 100, 300, 1000} & True & 25.3 \\
10 & \texttt{the auxiliary torus today (M8 = 10 TeV, M\_\allowbreak{}Pl = 2.435 10\textasciicircum{}18 GeV): V\_\allowbreak{}perp = M\_\allowbreak{}Pl\textasciicircum{}2/\allowbreak{}M8\textasciicircum{}6 = 5.929 10\textasciicircum{}12 GeV\textasciicircum{}-4, l\_\allowbreak{}perp = V\_\allowbreak{}perp\textasciicircum{}(1/\allowbreak{}4) = b0 Lc = 1560.4 GeV\textasciicircum{}-1 = 3.079 10\textasciicircum{}-13 m, lowest Kaluza-Klein mass 2 Pi/\allowbreak{}l\_\allowbreak{}perp = 4.027 MeV; N(mu) = V\_\allowbreak{}perp mu\textasciicircum{}4/\allowbreak{}(32 Pi\textasciicircum{}2) = 1.17 10\textasciicircum{}13 levels below m\_\allowbreak{}r/\allowbreak{}2 = 5 GeV (the lattice count of section 22.2 agrees to 10\textasciicircum{}-5) and N(M8) = M\_\allowbreak{}Pl\textasciicircum{}2/\allowbreak{}(32 Pi\textasciicircum{}2 M8\textasciicircum{}2)} & True & 25.3 \\
11 & \texttt{Newton's law in eight dimensions: Omega\_\allowbreak{}6 = 2 Pi\textasciicircum{}(7/\allowbreak{}2)/\allowbreak{}Gamma(7/\allowbreak{}2) = 16 Pi\textasciicircum{}3/\allowbreak{}15, r\textasciicircum{}-5 is harmonic away from r = 0 in seven space dimensions ((1/\allowbreak{}r\textasciicircum{}6) d/\allowbreak{}dr (r\textasciicircum{}6 d/\allowbreak{}dr r\textasciicircum{}-5) = 0), and grad\textasciicircum{}2 Phi = 8 Pi G8 ((D - 3)/\allowbreak{}(D - 2)) M delta\textasciicircum{}7 (D = 8) gives Phi = -(5/\allowbreak{}6) kappa M/\allowbreak{}(5 Omega\_\allowbreak{}6 r\textasciicircum{}5) = -5 kappa M/\allowbreak{}(32 Pi\textasciicircum{}3 r\textasciicircum{}5); the same formula in D = 4 gives -G M/\allowbreak{}r} & True & 25.4 \\
12 & \texttt{images = Kaluza-Klein modes (Poisson summation; lengths in units of l\_\allowbreak{}perp): Sum\_\allowbreak{}n (r\textasciicircum{}2 + |n|\textasciicircum{}2)\textasciicircum{}(-5/\allowbreak{}2) == (4 Pi\textasciicircum{}2/\allowbreak{}(3 r)) Sum\_\allowbreak{}m Exp[-2 Pi r |m|] to 10\textasciicircum{}-8 at r = 1/\allowbreak{}4, 1/\allowbreak{}2, 1, 2 (left: shells with r\_\allowbreak{}4(k) to k = 1.5 10\textasciicircum{}6 and the continuum tail Pi\textasciicircum{}2 Int k (r\textasciicircum{}2 + k)\textasciicircum{}(-5/\allowbreak{}2) dk; right: shells to k = 20000)} & True & 25.4 \\
13 & \texttt{near and far field: Z(5/\allowbreak{}2) = Sum\_\allowbreak{}(n != 0) |n|\textasciicircum{}-5 = Sum\_\allowbreak{}k r\_\allowbreak{}4(k) k\textasciicircum{}(-5/\allowbreak{}2) = 7 zeta(5/\allowbreak{}2) zeta(3/\allowbreak{}2) = 24.531 (shell sum with tail, to 10\textasciicircum{}-8); near the source r\textasciicircum{}5 S(r) - 1 -> r\textasciicircum{}5 Z(5/\allowbreak{}2) (the ratio lies between 0.98 and 1 at r = 1/\allowbreak{}10 and 1/\allowbreak{}20, closer to 1 at 1/\allowbreak{}20): the eight-dimensional law; far from it 3 r S/\allowbreak{}(4 Pi\textasciicircum{}2) - 1 -> 8 Exp[-2 Pi r], the 8 lightest Kaluza-Klein modes (ratio 1 to 2 10\textasciicircum{}-2 at r = 2 and 3)} & True & 25.4 \\
14 & \texttt{graviton + radion: far from the source the images add up to (5 kappa M/\allowbreak{}(32 Pi\textasciicircum{}3)) (4 Pi\textasciicircum{}2/\allowbreak{}(3 V\_\allowbreak{}perp r)) == (5/\allowbreak{}3) G4 M/\allowbreak{}r with G4 = kappa/\allowbreak{}(8 Pi V\_\allowbreak{}perp); the radion of section 20.3, coupled to matter by (phi/\allowbreak{}(Sqrt[3] M\_\allowbreak{}Pl)) T, gives the potential -(1/\allowbreak{}(Sqrt[3] M\_\allowbreak{}Pl))\textasciicircum{}2 M M' Exp[-m\_\allowbreak{}r r]/\allowbreak{}(4 Pi r), and (1/\allowbreak{}(Sqrt[3] M\_\allowbreak{}Pl))\textasciicircum{}2/\allowbreak{}(4 Pi) == (2/\allowbreak{}3) G4 (8 Pi G4 = 1/\allowbreak{}M\_\allowbreak{}Pl\textasciicircum{}2): graviton 1 + massless radion 2/\allowbreak{}3 = 5/\allowbreak{}3} & True & 25.4 \\
15 & \texttt{Kaluza-Klein momentum is conserved in the bulk and not at the brane: Int b\textasciicircum{}4 d\textasciicircum{}4y phi\_\allowbreak{}n phi\_\allowbreak{}m phi\_\allowbreak{}l\textasciicircum{}* == 1/\allowbreak{}Sqrt[V\_\allowbreak{}perp] for l = n + m and 0 otherwise (exact: n = (1, 0, 0, 2), m = (0, 1, -1, 0), l = (1, 1, -1, 2) and (1, 1, -1, 1)); at y = 0, phi\_\allowbreak{}n phi\_\allowbreak{}m phi\_\allowbreak{}l\textasciicircum{}* == V\_\allowbreak{}perp\textasciicircum{}(-3/\allowbreak{}2) for every n, m, l} & True & 25.5 \\
16 & \texttt{the delta function of the torus: Int\_\allowbreak{}(-Lc/\allowbreak{}2)\textasciicircum{}(Lc/\allowbreak{}2) Exp[-I q\_\allowbreak{}n y] delta(y) dy == 1 for every integer n, so delta\textasciicircum{}4(y) = Lc\textasciicircum{}-4 Sum\_\allowbreak{}n Exp[I q\_\allowbreak{}n.y] (every mode with the weight 1/\allowbreak{}Lc\textasciicircum{}4); a brane source rho4 delta\textasciicircum{}4(y)/\allowbreak{}b\textasciicircum{}4 has Int b\textasciicircum{}4 d\textasciicircum{}4y (rho4 delta\textasciicircum{}4(y)/\allowbreak{}b\textasciicircum{}4) == rho4, i.e. the average rho4/\allowbreak{}V\_\allowbreak{}perp over the torus} & True & 25.5 \\
17 & \texttt{the world brane W = \{x4 = x5 = x6 = x7 = 0\}: the induced metric h = J\textasciicircum{}T.g.J of the embedding X = (x0, x1, x2, x3, 0, 0, 0, 0) is, for the cosmological metric of section 11.2, the Friedmann-Robertson-Walker metric diag(1, -a\textasciicircum{}2, -a\textasciicircum{}2, -a\textasciicircum{}2) (independent of b); for the inhomogeneous metric of section 11.3 it is diag(n\textasciicircum{}2, -a\textasciicircum{}2, -a\textasciicircum{}2, -a\textasciicircum{}2) at x4 = 0 with n = 1 + x0\textasciicircum{}2/\allowbreak{}11, a = 1 + x0/\allowbreak{}3} & True & 25.6 \\
18 & \texttt{brane integrals in eight-dimensional form: at the brane Sqrt|det g| == Sqrt|det h| b\textasciicircum{}4 with b\textasciicircum{}4 = Sqrt[det(-g\_\allowbreak{}kl)] (cosmological metric; inhomogeneous metric exactly at x0 = 1/\allowbreak{}2), so Int\_\allowbreak{}W d\textasciicircum{}4x Sqrt|h| L == Int d\textasciicircum{}8X Sqrt|g| (delta\textasciicircum{}4(y)/\allowbreak{}b\textasciicircum{}4) L; test: Int d\textasciicircum{}4y b\textasciicircum{}4 (delta\textasciicircum{}4(y)/\allowbreak{}b\textasciicircum{}4) F(y) == F(0) = 1 for F = Exp[x4 + 2 x5] Cos[x6 - x7]} & True & 25.6 \\
19 & \texttt{how the brane sits in the bulk: the Christoffel symbols \{k mu nu\} at the brane (k = 4..7, mu, nu = 0..3; = (1/\allowbreak{}(2 b\textasciicircum{}2)) d\_\allowbreak{}k g\_\allowbreak{}mu nu for these metrics, b times the extrinsic curvature) all vanish for the cosmological metric: the brane is totally geodesic, and a free particle moving along it stays on it. For the inhomogeneous metric at x0 = 1/\allowbreak{}2 the only non-zero one is \{4 00\} = n (d\_\allowbreak{}4 n)/\allowbreak{}b\textasciicircum{}2 = 1458/\allowbreak{}58619: a free particle at rest on the brane is accelerated off it} & True & 25.6 \\
20 & \texttt{explicit forms: every Gamma17[A] is a signed permutation matrix (one entry +-1 in each row and column), and 2 Sqrt[2] U has only the entries 0, +-1, +-I, eight non-zero ones in each row and in each column} & True & 25.7 \\
21 & \texttt{brane spinor (x) transverse spinor: Gamma17[mu].Gamma17[nu] == U.(gamma\_\allowbreak{}mu.gamma\_\allowbreak{}nu (x) 1\_\allowbreak{}4).U\textasciicircum{}dag, Gamma17[k].Gamma17[l] == U.(1\_\allowbreak{}4 (x) rho\_\allowbreak{}k.rho\_\allowbreak{}l).U\textasciicircum{}dag and Gamma17[mu].Gamma17[k] == U.(gamma\_\allowbreak{}mu.gamma5 (x) rho\_\allowbreak{}k).U\textasciicircum{}dag (mu, nu = 0..3, k, l = 4..7): the Lorentz group of the brane acts on the brane spinor index, the rotations about the brane on the transverse index, and the generators that mix them are broken by the brane} & True & 25.7 \\
22 & \texttt{the transverse index carries Spin(4) = SU(2) x SU(2): rho4567 == -gamma5 (4 x 4), it commutes with every rho\_\allowbreak{}k.rho\_\allowbreak{}l and anticommutes with every rho\_\allowbreak{}k; for each even permutation (k, l, m, n) of (4, 5, 6, 7) the combination rho\_\allowbreak{}k.rho\_\allowbreak{}l + s rho\_\allowbreak{}m.rho\_\allowbreak{}n vanishes on the eigenspace rho4567 = s, and rho\_\allowbreak{}k.rho\_\allowbreak{}l - s rho\_\allowbreak{}m.rho\_\allowbreak{}n acts there with rank 2: the four values of a form (2, 1) + (1, 2)} & True & 25.7 \\
23 & \texttt{chirality at the brane: PL17 == U.(P\_\allowbreak{}L (x) P\_\allowbreak{}+ + P\_\allowbreak{}R (x) P\_\allowbreak{}-).U\textasciicircum{}dag with P\_\allowbreak{}L,R = (1 -+ gamma5)/\allowbreak{}2 and P\_\allowbreak{}+- = (1 +- rho4567)/\allowbreak{}2: the left-handed half psi\_\allowbreak{}L (chi17 = -1, 8 components) is two brane-left-handed Weyl spinors with rho4567 = +1 and two brane-right-handed ones with rho4567 = -1 (each part of rank 4)} & True & 25.7 \\
24 & \texttt{an explicit spinor: psi = U.(u (x) e\_\allowbreak{}1) with the brane spinor u = (1, 2, 0, -1): psibar psi == ubar u == 4 and psibar Gamma\textasciicircum{}mu psi == ubar gamma\textasciicircum{}mu u (mu = 0..3)} & True & 25.7 \\
25 & \texttt{brane tension: for S\_\allowbreak{}W = -Int d\textasciicircum{}4x Sqrt|h| f\_\allowbreak{}W and h = diag(1/\allowbreak{}u0, -1/\allowbreak{}u1, -1/\allowbreak{}u2, -1/\allowbreak{}u3) (h\textasciicircum{}00 = u0, h\textasciicircum{}ii = -u\_\allowbreak{}i), T\_\allowbreak{}mu nu = (2/\allowbreak{}Sqrt|h|) dS\_\allowbreak{}W/\allowbreak{}dh\textasciicircum{}mu nu == f\_\allowbreak{}W h\_\allowbreak{}mu nu, T\textasciicircum{}0\_\allowbreak{}0 == f\_\allowbreak{}W (energy density) and T\textasciicircum{}1\_\allowbreak{}1 == f\_\allowbreak{}W (pressure -f\_\allowbreak{}W: w = -1 along the brane)} & True & 25.8 \\
26 & \texttt{sources smeared over the torus (W(a, b) the comoving energy, section 20.2): the brane tension W = a\textasciicircum{}3 f\_\allowbreak{}W/\allowbreak{}Lc\textasciicircum{}4 gives eps = f\_\allowbreak{}W/\allowbreak{}V\_\allowbreak{}perp, pa = -eps, pb = 0; brane radiation W = a\textasciicircum{}3 r0 a\textasciicircum{}-4/\allowbreak{}Lc\textasciicircum{}4 gives pa = eps/\allowbreak{}3, pb = 0; radiation of the bulk, isotropic in the seven space directions, has <cos\textasciicircum{}2 theta> = 3/\allowbreak{}7 and <sin\textasciicircum{}2 theta> = 4/\allowbreak{}7 for the weight cos\textasciicircum{}2 theta sin\textasciicircum{}3 theta of the momentum sphere, so pa = (1/\allowbreak{}3)(3/\allowbreak{}7) eps = eps/\allowbreak{}7 and pb = (1/\allowbreak{}4)(4/\allowbreak{}7) eps = eps/\allowbreak{}7} & True & 25.8 \\
27 & \texttt{the brane in four-dimensional physics: (i) the tension is vacuum energy, kappa f\_\allowbreak{}W/\allowbreak{}V\_\allowbreak{}perp == f\_\allowbreak{}W/\allowbreak{}M\_\allowbreak{}Pl\textasciicircum{}2 (V\_\allowbreak{}perp = kappa M\_\allowbreak{}Pl\textasciicircum{}2), and f\_\allowbreak{}W = M8\textasciicircum{}4 is between 3 10\textasciicircum{}62 and 5 10\textasciicircum{}62 times rho\_\allowbreak{}Lambda = 3 H0\textasciicircum{}2 M\_\allowbreak{}Pl\textasciicircum{}2 Omega\_\allowbreak{}Lambda (section 24.5); (ii) normalcy temperature: (T\textasciicircum{}7/\allowbreak{}M8\textasciicircum{}6)/\allowbreak{}(T\textasciicircum{}2/\allowbreak{}M\_\allowbreak{}Pl) == 1 at T* = (M8\textasciicircum{}6/\allowbreak{}M\_\allowbreak{}Pl)\textasciicircum{}(1/\allowbreak{}5) = 13.3 GeV, and V\_\allowbreak{}perp T\textasciicircum{}4 T\textasciicircum{}3/\allowbreak{}M\_\allowbreak{}Pl\textasciicircum{}2 == T\textasciicircum{}7/\allowbreak{}M8\textasciicircum{}6; (iii) thin brane: (1/\allowbreak{}y\textasciicircum{}3) d/\allowbreak{}dy (y\textasciicircum{}3 d/\allowbreak{}dy y\textasciicircum{}-2) == 0 and 2 Pi\textasciicircum{}2 y\textasciicircum{}3 d/\allowbreak{}dy y\textasciicircum{}-2 == -4 Pi\textasciicircum{}2 (the Green's function -1/\allowbreak{}(4 Pi\textasciicircum{}2 |y|\textasciicircum{}2) of the four transverse directions), and Sqrt[kappa f\_\allowbreak{}W] == 1/\allowbreak{}M8 for f\_\allowbreak{}W = M8\textasciicircum{}4} & True & 25.8 \\*
\end{longtable}}
The comparison with the 4-sphere (section 10 of the document) rests on section 26 of 7+1.nb: 33 checks, 33 True; the independent program tools/ref\_radion.py repeats them with 31 boolean checks, 31 True, and the 15 groups of numbers printed by both agree to a largest relative difference of $8.2\times10^{-11}$.

\Needspace{8\baselineskip}\subsection*{6.2 The independent Python program}\addcontentsline{toc}{subsection}{6.2 The independent Python program}
tools/ref\_torus\_brane.py — its boolean checks:

{\footnotesize
\begin{longtable}{>{\raggedright\arraybackslash}p{\dimexpr 0.0550\linewidth-2\tabcolsep\relax}>{\raggedright\arraybackslash}p{\dimexpr 0.8250\linewidth-2\tabcolsep\relax}>{\raggedright\arraybackslash}p{\dimexpr 0.1200\linewidth-2\tabcolsep\relax}}
\toprule
\textbf{\#} & \textbf{check} & \textbf{result} \\
\midrule\endfirsthead
\toprule
\textbf{\#} & \textbf{check} & \textbf{result} \\
\midrule\endhead
\endfoot
\bottomrule\endlastfoot
1 & \texttt{torus: det(-g\_\allowbreak{}kl) = b\textasciicircum{}8 for the cosmological metric; the proper volume of [0, Lc)\textasciicircum{}4 is b\textasciicircum{}4 Lc\textasciicircum{}4} & True \\
2 & \texttt{only b Lc is physical: y = lambda y' maps (Lc, b) to (Lc/\allowbreak{}lambda, lambda b); b\textasciicircum{}2 dy.dy, b Lc and b\textasciicircum{}4 Lc\textasciicircum{}4 unchanged} & True \\
3 & \texttt{symmetries fixing y = 0: 384 = 2\textasciicircum{}4 4! distinct signed permutation matrices, orthogonal integer matrices; shells 1, 2, 3, 4, 9 have orbits \{8\}, \{24\}, \{32\}, \{8, 16\}, \{8, 96\}} & True \\
4 & \texttt{Fourier modes phi\_\allowbreak{}n = exp(i q\_\allowbreak{}n.y)/\allowbreak{}sqrt(V\_\allowbreak{}perp), q\_\allowbreak{}n = 2 pi n/\allowbreak{}Lc: orthonormal for b\textasciicircum{}4 d\textasciicircum{}4y (three exact pairs)} & True \\
5 & \texttt{Parseval: f = prod cos\textasciicircum{}2(pi y\_\allowbreak{}k/\allowbreak{}Lc) has 81 coefficients sqrt(V\_\allowbreak{}perp) prod w(n\_\allowbreak{}k) (3 checked by integration), Sum |c\_\allowbreak{}n|\textasciicircum{}2 = Int b\textasciicircum{}4 f\textasciicircum{}2 = (3/\allowbreak{}8)\textasciicircum{}4 V\_\allowbreak{}perp} & True \\
6 & \texttt{Laplacian of the torus: -(1/\allowbreak{}b\textasciicircum{}2) Sum d\textasciicircum{}2/\allowbreak{}dy\_\allowbreak{}k\textasciicircum{}2 phi\_\allowbreak{}n = (|q\_\allowbreak{}n|\textasciicircum{}2/\allowbreak{}b\textasciicircum{}2) phi\_\allowbreak{}n} & True \\
7 & \texttt{scalar Kaluza-Klein masses: (Box + m\textasciicircum{}2) exp(-iEx0 + ik.x + iq.y) = (-E\textasciicircum{}2 + |k|\textasciicircum{}2 + |q|\textasciicircum{}2/\allowbreak{}b0\textasciicircum{}2 + m\textasciicircum{}2) exp(...) on diag(1,-1,-1,-1,-b0\textasciicircum{}2,...)} & True \\
8 & \texttt{graviton on the torus: D(D-3)/\allowbreak{}2 = 20 (D = 8) = 2 + 8 + 10 (zero mode: graviton, 4 graviphotons, 10 moduli) = 5 + 9 + 6 (massive level)} & True \\
9 & \texttt{Jacobi: \#\{n in Z\textasciicircum{}4 : |n|\textasciicircum{}2 = k\} = r\_\allowbreak{}4(k) = 8 sigma(k) - 32 sigma(k/\allowbreak{}4), direct count for k = 1..100} & True \\
10 & \texttt{Weyl's law: N(10) = 49689, N(100) = 493490641; |N(R)/\allowbreak{}(pi\textasciicircum{}2 R\textasciicircum{}4/\allowbreak{}2) - 1| < 1e-2, 1e-3, 3e-5, 5e-6, 1e-6} & True \\
11 & \texttt{numbers today: V\_\allowbreak{}perp = 5.929e12 GeV\textasciicircum{}-4, l\_\allowbreak{}perp = 1560.4 GeV\textasciicircum{}-1 = 3.079e-13 m, 2 pi/\allowbreak{}l\_\allowbreak{}perp = 4.027 MeV; N(5 GeV) = 1.17e13 (lattice to 1e-5); N(M8) = M\_\allowbreak{}Pl\textasciicircum{}2/\allowbreak{}(32 pi\textasciicircum{}2 M8\textasciicircum{}2)} & True \\
12 & \texttt{from eight to four dimensions: for diag(1, -a\textasciicircum{}2, -a\textasciicircum{}2, -a\textasciicircum{}2, -b0\textasciicircum{}2, ..., -b0\textasciicircum{}2) (static torus) R = R4 = -6(Ha' + 2 Ha\textasciicircum{}2), the scalar curvature of the FRW metric diag(1, -a\textasciicircum{}2, -a\textasciicircum{}2, -a\textasciicircum{}2); sqrt|g| = a\textasciicircum{}3 b0\textasciicircum{}4} & True \\
13 & \texttt{quaternion subalgebra of the brane directions: e1 e2 = -e3, e2 e3 = -e1, e3 e1 = -e2 (e\_\allowbreak{}i e\_\allowbreak{}j = -delta\_\allowbreak{}ij - eps\_\allowbreak{}ijk e\_\allowbreak{}k), span\{e0, ..., e3\} closed under the octonion product, Re(e\_\allowbreak{}A e\_\allowbreak{}B) = diag(1, -1, -1, -1)} & True \\
14 & \texttt{Newton in 8D: Omega\_\allowbreak{}6 = 16 pi\textasciicircum{}3/\allowbreak{}15, (1/\allowbreak{}r\textasciicircum{}6) d/\allowbreak{}dr (r\textasciicircum{}6 d/\allowbreak{}dr r\textasciicircum{}-5) = 0, Phi = -(5/\allowbreak{}6) kappa M/\allowbreak{}(5 Omega\_\allowbreak{}6 r\textasciicircum{}5) = -5 kappa M/\allowbreak{}(32 pi\textasciicircum{}3 r\textasciicircum{}5); D = 4: -G M/\allowbreak{}r} & True \\
15 & \texttt{4D Fourier transform: Int d\textasciicircum{}4y exp(-ip.y) (r\textasciicircum{}2 + |y|\textasciicircum{}2)\textasciicircum{}(-5/\allowbreak{}2) = (4 pi\textasciicircum{}2/\allowbreak{}(3 r)) exp(-p r) (mpmath, three (p, r), to 1e-20)} & True \\
16 & \texttt{Poisson duality: Sum\_\allowbreak{}n (r\textasciicircum{}2 + |n|\textasciicircum{}2)\textasciicircum{}(-5/\allowbreak{}2) = (4 pi\textasciicircum{}2/\allowbreak{}(3 r)) Sum\_\allowbreak{}m exp(-2 pi r |m|) to 1e-8 at r = 1/\allowbreak{}4, 1/\allowbreak{}2, 1, 2} & True \\
17 & \texttt{Z(5/\allowbreak{}2) = 7 zeta(5/\allowbreak{}2) zeta(3/\allowbreak{}2) = Sum r\_\allowbreak{}4(k) k\textasciicircum{}-5/\allowbreak{}2 (to 1e-8); near field ratio -> 1 (0.98 < . < 1 at r = 1/\allowbreak{}10, 1/\allowbreak{}20); far field ratio -> 1 (to 2e-2 at r = 2, 3)} & True \\
18 & \texttt{graviton + radion: (5 kappa M/\allowbreak{}(32 pi\textasciicircum{}3))(4 pi\textasciicircum{}2/\allowbreak{}(3 V\_\allowbreak{}perp r)) = (5/\allowbreak{}3) G4 M/\allowbreak{}r; radion (1/\allowbreak{}(sqrt3 M\_\allowbreak{}Pl))\textasciicircum{}2/\allowbreak{}(4 pi) = (2/\allowbreak{}3)/\allowbreak{}(8 pi M\_\allowbreak{}Pl\textasciicircum{}2)} & True \\
19 & \texttt{Kaluza-Klein momentum: bulk vertex Int b\textasciicircum{}4 phi\_\allowbreak{}n phi\_\allowbreak{}m phi\_\allowbreak{}l\textasciicircum{}* = 1/\allowbreak{}sqrt(V\_\allowbreak{}perp) for l = n + m, 0 otherwise; at y = 0 phi\_\allowbreak{}n phi\_\allowbreak{}m phi\_\allowbreak{}l\textasciicircum{}* = V\_\allowbreak{}perp\textasciicircum{}(-3/\allowbreak{}2)} & True \\
20 & \texttt{delta function of the torus: Fourier coefficients Int exp(-i q\_\allowbreak{}n y) delta(y) dy = 1; Int b\textasciicircum{}4 d\textasciicircum{}4y rho4 delta\textasciicircum{}4(y)/\allowbreak{}b\textasciicircum{}4 = rho4 (average rho4/\allowbreak{}V\_\allowbreak{}perp)} & True \\
21 & \texttt{induced metric h = J\textasciicircum{}T g J: FRW diag(1, -a\textasciicircum{}2, -a\textasciicircum{}2, -a\textasciicircum{}2) for the cosmological metric; diag(n\textasciicircum{}2, -a\textasciicircum{}2, -a\textasciicircum{}2, -a\textasciicircum{}2) at x4 = 0 with n = 1 + x0\textasciicircum{}2/\allowbreak{}11, a = 1 + x0/\allowbreak{}3 for the inhomogeneous metric} & True \\
22 & \texttt{brane volume element: sqrt|g| = sqrt|h| b\textasciicircum{}4 at the brane (cosmological metric; inhomogeneous metric at x0 = 1/\allowbreak{}2, exact), so Int\_\allowbreak{}W d\textasciicircum{}4x sqrt|h| L = Int d\textasciicircum{}8X sqrt|g| delta\textasciicircum{}4(y)/\allowbreak{}b\textasciicircum{}4 L; test Int b\textasciicircum{}4 (delta\textasciicircum{}4/\allowbreak{}b\textasciicircum{}4) F = F(0)} & True \\
23 & \texttt{Christoffel symbols \{k mu nu\} at the brane: all 0 for the cosmological metric (totally geodesic); for the inhomogeneous metric at x0 = 1/\allowbreak{}2 only \{4 00\} = 1458/\allowbreak{}58619 = n d4n/\allowbreak{}b\textasciicircum{}2} & True \\
24 & \texttt{R\_\allowbreak{}A = (gamma\_\allowbreak{}mu (x) 1\_\allowbreak{}4, gamma5 (x) rho\_\allowbreak{}k) and Gamma17[A] both satisfy X\_\allowbreak{}A X\_\allowbreak{}B + X\_\allowbreak{}B X\_\allowbreak{}A = 2 eta\_\allowbreak{}AB 1\_\allowbreak{}16} & True \\
25 & \texttt{every Gamma17[A] is a signed permutation matrix} & True \\
26 & \texttt{intertwiner: Schur average with the seed matrix unit (1, 2); 2 sqrt(2) U has entries 0, +-1, +-i (rounding error < 1e-12) with 8 non-zero entries in each row and column; Gamma17[A].U8 = U8.R\_\allowbreak{}A and U8\textasciicircum{}dag.U8 = 8 1\_\allowbreak{}16 exactly} & True \\
27 & \texttt{Schur's lemma: the linear map X -> (Gamma\_\allowbreak{}A X - X R\_\allowbreak{}A)\_\allowbreak{}A has a one-dimensional null space (one singular value < 1e-10, the next > 0.1)} & True \\
28 & \texttt{brane Lorentz group on the brane spinor index, transverse rotations on a: Gamma\_\allowbreak{}mu Gamma\_\allowbreak{}nu = U (gamma\_\allowbreak{}mu gamma\_\allowbreak{}nu (x) 1) U\textasciicircum{}dag, Gamma\_\allowbreak{}k Gamma\_\allowbreak{}l = U (1 (x) rho\_\allowbreak{}k rho\_\allowbreak{}l) U\textasciicircum{}dag, Gamma\_\allowbreak{}mu Gamma\_\allowbreak{}k = U (gamma\_\allowbreak{}mu gamma5 (x) rho\_\allowbreak{}k) U\textasciicircum{}dag} & True \\
29 & \texttt{Spin(4) on the transverse index: rho4567 = -gamma5 (4x4) commutes with rho\_\allowbreak{}k rho\_\allowbreak{}l, anticommutes with rho\_\allowbreak{}k; (rho\_\allowbreak{}k rho\_\allowbreak{}l + s rho\_\allowbreak{}m rho\_\allowbreak{}n) P\_\allowbreak{}s = 0 and (rho\_\allowbreak{}k rho\_\allowbreak{}l - s rho\_\allowbreak{}m rho\_\allowbreak{}n) P\_\allowbreak{}s has rank 2 for even (k,l,m,n): 4 = (2,1) + (1,2)} & True \\
30 & \texttt{chirality: chi = i Gamma\_\allowbreak{}0...Gamma\_\allowbreak{}7 = U (gamma5 (x) rho4567) U\textasciicircum{}dag; P\_\allowbreak{}L17 = U (P\_\allowbreak{}L (x) P\_\allowbreak{}+ + P\_\allowbreak{}R (x) P\_\allowbreak{}-) U\textasciicircum{}dag, each part of rank 4} & True \\
31 & \texttt{example psi = U.(u (x) e\_\allowbreak{}1), u = (1, 2, 0, -1): psibar psi = ubar u = 4 and psibar Gamma\textasciicircum{}mu psi = ubar gamma\textasciicircum{}mu u} & True \\
32 & \texttt{four brane Dirac spinors: psibar psi' = Sum\_\allowbreak{}a psibar\_\allowbreak{}a psi'\_\allowbreak{}a and psibar Gamma\textasciicircum{}mu psi' = Sum\_\allowbreak{}a psibar\_\allowbreak{}a gamma\textasciicircum{}mu psi'\_\allowbreak{}a (random integer spinors)} & True \\
33 & \texttt{Kaluza-Klein operator X = gamma5 (x) (rho.q)/\allowbreak{}b0: X\textasciicircum{}2 = -|q|\textasciicircum{}2/\allowbreak{}b0\textasciicircum{}2 1\_\allowbreak{}16 and X anticommutes with gamma\_\allowbreak{}mu (x) 1\_\allowbreak{}4} & True \\
34 & \texttt{brane tension: T\_\allowbreak{}mu nu = (2/\allowbreak{}sqrt|h|) dS/\allowbreak{}dh\textasciicircum{}mu nu = f\_\allowbreak{}W h\_\allowbreak{}mu nu for S = -Int sqrt|h| f\_\allowbreak{}W; T\textasciicircum{}0\_\allowbreak{}0 = f\_\allowbreak{}W, T\textasciicircum{}1\_\allowbreak{}1 = f\_\allowbreak{}W (p = -f\_\allowbreak{}W, w = -1)} & True \\
35 & \texttt{smeared sources: tension eps = f\_\allowbreak{}W/\allowbreak{}V\_\allowbreak{}perp, pa = -eps, pb = 0; brane radiation pa = eps/\allowbreak{}3, pb = 0; bulk radiation <cos\textasciicircum{}2> = 3/\allowbreak{}7, <sin\textasciicircum{}2> = 4/\allowbreak{}7: pa = pb = eps/\allowbreak{}7} & True \\
36 & \texttt{normalcy temperature: (T\textasciicircum{}7/\allowbreak{}M8\textasciicircum{}6)/\allowbreak{}(T\textasciicircum{}2/\allowbreak{}M\_\allowbreak{}Pl) = 1 at T* = (M8\textasciicircum{}6/\allowbreak{}M\_\allowbreak{}Pl)\textasciicircum{}(1/\allowbreak{}5) = 13.3 GeV; V\_\allowbreak{}perp T\textasciicircum{}4 T\textasciicircum{}3/\allowbreak{}M\_\allowbreak{}Pl\textasciicircum{}2 = T\textasciicircum{}7/\allowbreak{}M8\textasciicircum{}6} & True \\
37 & \texttt{the tension is vacuum energy: kappa f\_\allowbreak{}W/\allowbreak{}V\_\allowbreak{}perp = f\_\allowbreak{}W/\allowbreak{}M\_\allowbreak{}Pl\textasciicircum{}2 (V\_\allowbreak{}perp = kappa M\_\allowbreak{}Pl\textasciicircum{}2); f\_\allowbreak{}W = M8\textasciicircum{}4 is 3e62 to 5e62 rho\_\allowbreak{}Lambda} & True \\
38 & \texttt{Green's function of the four transverse directions: (1/\allowbreak{}y\textasciicircum{}3) d/\allowbreak{}dy (y\textasciicircum{}3 d/\allowbreak{}dy y\textasciicircum{}-2) = 0, flux 2 pi\textasciicircum{}2 y\textasciicircum{}3 d/\allowbreak{}dy y\textasciicircum{}-2 = -4 pi\textasciicircum{}2; sqrt(kappa f\_\allowbreak{}W) = 1/\allowbreak{}M8 for f\_\allowbreak{}W = M8\textasciicircum{}4} & True \\*
\end{longtable}}
\textbf{Numbers printed by both programs} (largest relative difference per group, and the tolerance used by tools/verify\_7p1\_nb.py; integers and exact matrices must agree exactly):

{\footnotesize
\begin{longtable}{>{\raggedright\arraybackslash}p{\dimexpr 0.5200\linewidth-2\tabcolsep\relax}>{\raggedright\arraybackslash}p{\dimexpr 0.1000\linewidth-2\tabcolsep\relax}>{\raggedright\arraybackslash}p{\dimexpr 0.2200\linewidth-2\tabcolsep\relax}>{\raggedright\arraybackslash}p{\dimexpr 0.1600\linewidth-2\tabcolsep\relax}}
\toprule
\textbf{quantity} & \textbf{values} & \textbf{max. relative difference} & \textbf{tolerance} \\
\midrule\endfirsthead
\toprule
\textbf{quantity} & \textbf{values} & \textbf{max. relative difference} & \textbf{tolerance} \\
\midrule\endhead
\endfoot
\bottomrule\endlastfoot
orbit sizes of the shells $k=1,\dots,12$ & 16 & 0 & exact \\
$r_4(k)$, $k=0,\dots,16$ & 17 & 0 & exact \\
Weyl: $N(R)$, $R=10,30,100,300,1000$ & 5 & 0 & exact \\
Weyl: $N(R)/(\pi^2R^4/2)-1$ & 5 & $3.8\times10^{-10}$ & $1\times10^{-9}$ \\
$V_\perp$, $\ell_\perp$ (GeV$^{-1}$, m, fm), $2\pi/\ell_\perp$ & 5 & $2.2\times10^{-16}$ & $1\times10^{-12}$ \\
$\ell_\perp/r_p$, $a_0/\ell_\perp$, $R=\ell_\perp/2\pi$, $M_D$ & 5 & $1.7\times10^{-16}$ & $1\times10^{-12}$ \\
$N(\mu)$ at 5 GeV, $m_h$, $M_8$; lattice count below 5 GeV & 4 & $1.8\times10^{-16}$ & $1\times10^{-12}$ \\
the first levels ($k$, $r_4$, $4r_4$, mass) & 32 & $2.5\times10^{-16}$ & $1\times10^{-12}$ \\
image sums $S(r)$ and Kaluza–Klein sums, $r=\frac14,\frac12,1,2$ & 8 & $1.2\times10^{-15}$ & $1\times10^{-9}$ \\
$Z(5/2)$: closed form, shell sum & 2 & $2.9\times10^{-16}$ & $1\times10^{-9}$ \\
near-field and far-field ratios & 5 & $7.8\times10^{-12}$ & $1\times10^{-6}$ \\
Christoffel symbol $\{4\,00\}$ at the brane & 1 & 0 & $1\times10^{-15}$ \\
$\Gamma_A$ as signed permutations & 128 & 0 & exact \\
$\mathrm{Re}(2\sqrt2\,U)$ & 256 & 0 & exact \\
$\mathrm{Im}(2\sqrt2\,U)$ & 256 & 0 & exact \\
example spinor: $\bar\psi\psi$, $\bar\psi\Gamma^\mu\psi$ & 5 & 0 & $1\times10^{-15}$ \\
$T_*$, $f_W/\rho_\Lambda$, $1/(2\pi M_8)$ & 3 & $5.8\times10^{-16}$ & $1\times10^{-12}$ \\*
\end{longtable}}
\Needspace{8\baselineskip}\subsection*{6.3 The earlier sections are unchanged}\addcontentsline{toc}{subsection}{6.3 The earlier sections are unchanged}
tools/verify\_7p1\_nb.py compares the evaluated notebook with the backup of section 3 and with the Python program; its lines on section 25:

\begin{Verbatim}[fontsize=\small,frame=single,framesep=2mm]
section 25: check outputs (all True) : 27 (True)
stage 6: without section 25 (and with the old title of the summary section) the 786 input and
    ...text cells equal those of the backup taken before section 25 was appended : True
stage 6: the 481 outputs of sections 0-24 and of the summary equal those of the backup (summary
    ...counts and tables, the private context name and measured times excepted) : True
section 25 numbers: 17 groups of printed numbers agree with tools/ref_torus_brane.py (integers
    ...and exact values exactly, floating-point numbers to 1e-15 ... 1e-6 per group) : True
VERDICT: ALL CROSS-CHECKS PASSED
\end{Verbatim}
\Needspace{8\baselineskip}\subsection*{6.4 Files}\addcontentsline{toc}{subsection}{6.4 Files}
\textbf{SHA-256 of the files} (at the time this document was generated):

\begin{Verbatim}[fontsize=\small,frame=single,framesep=2mm]
7+1.nb                             41d252f9852f210c228f45e1ac956d2d7ef26faeb0062a5880280e1aae0ef
    ...fc7
tools/make_7p1_nb.py               996d26c9552cd44902560af16a1683fa71457261f18fefaeddae87bd99a2b
    ...144
tools/nb_part_gravity.py           c4caf2ad5c932e8f603341edade5f4b98fa784feb05906c051eab7c884cd6
    ...8f8
tools/nb_part_ks_gravity.py        a656c0ad040b4576b75b17e029545109e531c8e656ad99bfdb5cba946e967
    ...772
tools/nb_part_sm.py                43e92beed57453a4b57931de6d603ac6c218efba6e4005133b69b8054a975
    ...c5b
tools/nb_part_coupled.py           e630521d64132e71afc861114f699eb06454557cf14a57ad8ec1eef96247a
    ...ae0
tools/nb_part_dark.py              10ad8fb8ae25fa07bbee758c653995acd4e2c80733059a8b3648bd1593e46
    ...db0
tools/nb_part_torus_brane.py       fb904ff8d85b904c55b3b2a8f60fb0cb047f1ccf1b98db3364092f2bae0b2
    ...162
tools/nb_part_radion.py            53c5604707a3e22228c840ee0a2d2125616f9410f75a5535041dfaf18f497
    ...c00
tools/lovelock71.py                d93ecbd0f82adb972218d59dfe32999a4af944945ca9758607f3a633e1af9
    ...5a1
tools/ref_gravity.py               f07cf8adaff585af89c87d19532f1aaa7db606ad88ae90483b2d1f41341ed
    ...f36
tools/ref_ks_gravity.py            b6488ef24b1c16aacd8ff6cf6a5a1ba3f316e082449ad97480022d1a26383
    ...27f
tools/ref_sm.py                    f3f80806893f42bbf166e1f1de896ffd611dbfe847685da8e3ee367e9cde6
    ...aba
tools/ref_coupled.py               695dbd9375753653472d9c37502f0067c4a64fe01a64ba8c05e8b2c5da85a
    ...a58
tools/ref_dark.py                  1600e6b57ab0383b4f2b664e2796d3864c9afa77a9ed59f007f66e1139453
    ...ff5
tools/ref_torus_brane.py           1f623656ab0f5c8ee43cb310f7bfd580cbba5a7385da2963adaae4d2f8bda
    ...af7
tools/ref_radion.py                e9d47176b9f2d2232a84b3b82188bd3b0a19cde3dbe1ba013c1fe5abff634
    ...1ce
tools/ref71.py                     1a7a06eb0add3f459611776dc9c474a5bb959d947a9db949db4612b1ab943
    ...4e0
tools/ref44.py                     ab100c3ebcc055fbe59c132a02f80ecb17630e934e99d59f88e919ec8064f
    ...af5
\end{Verbatim}
\textbf{Output of tools/ref\_torus\_brane.py (independent checks of section 25; exit code 0)} (complete):

\begin{Verbatim}[fontsize=\small,frame=single,framesep=2mm]
torus: det(-g_kl) = b^8 for the cosmological metric; the proper volume of [0, Lc)^4 is b^4 Lc^4
    ...: True
only b Lc is physical: y = lambda y' maps (Lc, b) to (Lc/lambda, lambda b); b^2 dy.dy, b Lc and
    ...b^4 Lc^4 unchanged : True
symmetries fixing y = 0: 384 = 2^4 4! distinct signed permutation matrices, orthogonal integer
    ...matrices; shells 1, 2, 3, 4, 9 have orbits {8}, {24}, {32}, {8, 16}, {8, 96} : True
NUM orbit sizes of the shells k = 1..12 : [[8], [24], [32], [8, 16], [48], [96], [64], [24],
    ...[8, 96], [48, 96], [96], [32, 64]]
Fourier modes phi_n = exp(i q_n.y)/sqrt(V_perp), q_n = 2 pi n/Lc: orthonormal for b^4 d^4y
    ...(three exact pairs) : True
Parseval: f = prod cos^2(pi y_k/Lc) has 81 coefficients sqrt(V_perp) prod w(n_k) (3 checked by
    ...integration), Sum |c_n|^2 = Int b^4 f^2 = (3/8)^4 V_perp : True
Laplacian of the torus: -(1/b^2) Sum d^2/dy_k^2 phi_n = (|q_n|^2/b^2) phi_n : True
scalar Kaluza-Klein masses: (Box + m^2) exp(-iEx0 + ik.x + iq.y) = (-E^2 + |k|^2 + |q|^2/b0^2 +
    ...m^2) exp(...) on diag(1,-1,-1,-1,-b0^2,...) : True
graviton on the torus: D(D-3)/2 = 20 (D = 8) = 2 + 8 + 10 (zero mode: graviton, 4 graviphotons,
    ...10 moduli) = 5 + 9 + 6 (massive level) : True
Jacobi: #{n in Z^4 : |n|^2 = k} = r_4(k) = 8 sigma(k) - 32 sigma(k/4), direct count for k =
    ...1..100 : True
NUM r_4(k), k = 0..16 : [1, 8, 24, 32, 24, 48, 96, 64, 24, 104, 144, 96, 96, 112, 192, 192, 24]
NUM N(R) for R = 10, 30, 100, 300, 1000 : [49689, 4000425, 493490641, 39972070929,
    ...4934805110729]
NUM N(R)/(pi^2 R^4/2) - 1 for R = 10, 30, 100, 300, 1000 : [0.006909658800824392,
    ...0.0008093730182694436, 2.1117250719626668e-05, 4.330657224604906e-06,
    ...5.897266401877488e-07]
Weyl's law: N(10) = 49689, N(100) = 493490641; |N(R)/(pi^2 R^4/2) - 1| < 1e-2, 1e-3, 3e-5,
    ...5e-6, 1e-6 : True
NUM V_perp (GeV^-4), l_perp (GeV^-1), l_perp (m), l_perp (fm), 2 pi/l_perp (MeV) :
    ...(5929225000000.0, 1560.4486534327234, 3.079186208511254e-13, 307.9186208511254,
    ...4.026524867292262)
NUM l_perp/proton radius, Bohr radius/l_perp, R = l_perp/(2 pi) (GeV^-1, m), M_D = (2 pi)^(2/3)
    ...M8 (GeV) : (365.9598536381334, 171.8561902623129, 248.3531166349098,
    ...4.9006770578495764e-14, 34050.219214767545)
NUM N(mu) = V_perp mu^4/(32 pi^2) for mu = 5 GeV, m_h, M8; lattice count below 5 GeV :
    ...(11733517482065.238, 4.61280953636634e+18, 1.877362797130438e+26, 11733511217161)
numbers today: V_perp = 5.929e12 GeV^-4, l_perp = 1560.4 GeV^-1 = 3.079e-13 m, 2 pi/l_perp =
    ...4.027 MeV; N(5 GeV) = 1.17e13 (lattice to 1e-5); N(M8) = M_Pl^2/(32 pi^2 M8^2) : True
NUM first levels: k, r_4(k), Dirac fermions 4 r_4(k), mass 2 pi sqrt(k)/l_perp (MeV), k = 1..8
    ...: [(1, 8, 32, 4.026524867292262), (2, 24, 96, 5.694366076557245), (3, 32, 128,
    ...6.9741456480897295), (4, 24, 96, 8.053049734584524), (5, 48, 192, 9.003583316358819),
    ...(6, 96, 384, 9.862931361493793), (7, 64, 256, 10.65318344667268), (8, 24, 96,
    ...11.38873215311449)]
from eight to four dimensions: for diag(1, -a^2, -a^2, -a^2, -b0^2, ..., -b0^2) (static torus)
    ...R = R4 = -6(Ha' + 2 Ha^2), the scalar curvature of the FRW metric diag(1, -a^2, -a^2,
    ...-a^2); sqrt|g| = a^3 b0^4 : True
quaternion subalgebra of the brane directions: e1 e2 = -e3, e2 e3 = -e1, e3 e1 = -e2 (e_i e_j =
    ...-delta_ij - eps_ijk e_k), span{e0, ..., e3} closed under the octonion product, Re(e_A
    ...e_B) = diag(1, -1, -1, -1) : True
NUM quaternion products e1 e2, e2 e3, e3 e1: coefficients of e3, e1, e2 : (-1, -1, -1)
Newton in 8D: Omega_6 = 16 pi^3/15, (1/r^6) d/dr (r^6 d/dr r^-5) = 0, Phi = -(5/6) kappa M/(5
    ...Omega_6 r^5) = -5 kappa M/(32 pi^3 r^5); D = 4: -G M/r : True
4D Fourier transform: Int d^4y exp(-ip.y) (r^2 + |y|^2)^(-5/2) = (4 pi^2/(3 r)) exp(-p r)
    ...(mpmath, three (p, r), to 1e-20) : True
NUM images S(r), Kaluza-Klein (4 pi^2/(3 r)) Sum exp(-2 pi r |m|), relative difference, at r =
    ...1/4, 1/2, 1, 2 : [(1046.8493759696448, 1046.8493759704088, -7.297495940861154e-13),
    ...(51.12525179526186, 51.12525179602524, -1.4931500480486193e-11), (13.40960454140553,
    ...13.409604542168916, -5.692823989988938e-11), (6.579922931695463, 6.579922932458839,
    ...-1.1601597460497715e-10)]
Poisson duality: Sum_n (r^2 + |n|^2)^(-5/2) = (4 pi^2/(3 r)) Sum_m exp(-2 pi r |m|) to 1e-8 at
    ...r = 1/4, 1/2, 1, 2 : True
NUM Z(5/2) = 7 zeta(5/2) zeta(3/2), shell sum with tail : (24.53127768992603,
    ...24.531277689162653)
NUM near field (r^5 S - 1)/(r^5 Z(5/2)) at r = 1/20, 1/10; far field (3 r S/(4 pi^2) - 1)/(8
    ...exp(-2 pi r)) at r = 1, 2, 3 : [0.9970259928360821, 0.9882280058562034,
    ...1.272311659711708, 1.016880665840871, 1.0012238149051904]
Z(5/2) = 7 zeta(5/2) zeta(3/2) = Sum r_4(k) k^-5/2 (to 1e-8); near field ratio -> 1 (0.98 < . <
    ...1 at r = 1/10, 1/20); far field ratio -> 1 (to 2e-2 at r = 2, 3) : True
graviton + radion: (5 kappa M/(32 pi^3))(4 pi^2/(3 V_perp r)) = (5/3) G4 M/r; radion (1/(sqrt3
    ...M_Pl))^2/(4 pi) = (2/3)/(8 pi M_Pl^2) : True
Kaluza-Klein momentum: bulk vertex Int b^4 phi_n phi_m phi_l^* = 1/sqrt(V_perp) for l = n + m,
    ...0 otherwise; at y = 0 phi_n phi_m phi_l^* = V_perp^(-3/2) : True
delta function of the torus: Fourier coefficients Int exp(-i q_n y) delta(y) dy = 1; Int b^4
    ...d^4y rho4 delta^4(y)/b^4 = rho4 (average rho4/V_perp) : True
induced metric h = J^T g J: FRW diag(1, -a^2, -a^2, -a^2) for the cosmological metric;
    ...diag(n^2, -a^2, -a^2, -a^2) at x4 = 0 with n = 1 + x0^2/11, a = 1 + x0/3 for the
    ...inhomogeneous metric : True
brane volume element: sqrt|g| = sqrt|h| b^4 at the brane (cosmological metric; inhomogeneous
    ...metric at x0 = 1/2, exact), so Int_W d^4x sqrt|h| L = Int d^8X sqrt|g| delta^4(y)/b^4 L;
    ...test Int b^4 (delta^4/b^4) F = F(0) : True
Christoffel symbols {k mu nu} at the brane: all 0 for the cosmological metric (totally
    ...geodesic); for the inhomogeneous metric at x0 = 1/2 only {4 00} = 1458/58619 = n d4n/b^2
    ...: True
NUM Christoffel {4 00} at the brane (inhomogeneous metric, x0 = 1/2) : 0.02487248161858783
R_A = (gamma_mu (x) 1_4, gamma5 (x) rho_k) and Gamma17[A] both satisfy X_A X_B + X_B X_A = 2
    ...eta_AB 1_16 : True
every Gamma17[A] is a signed permutation matrix : True
NUM Gamma17[A] as signed permutations: column j has the entry sign(c) in row |c|, c listed for
    ...A = 0..7 : [[9, 10, 11, 12, 13, 14, 15, 16, 1, 2, 3, 4, 5, 6, 7, 8], [-8, -7, 6, 5, -4,
    ...-3, 2, 1, 16, 15, -14, -13, 12, 11, -10, -9], [7, -8, -5, 6, 3, -4, -1, 2, -15, 16, 13,
    ...-14, -11, 12, 9, -10], [-6, 5, -8, 7, -2, 1, -4, 3, 14, -13, 16, -15, 10, -9, 12, -11],
    ...[-2, 1, 4, -3, 6, -5, -8, 7, 10, -9, -12, 11, -14, 13, 16, -15], [3, 4, -1, -2, -7, -8,
    ...5, 6, -11, -12, 9, 10, 15, 16, -13, -14], [4, -3, 2, -1, -8, 7, -6, 5, -12, 11, -10, 9,
    ...16, -15, 14, -13], [5, 6, 7, 8, -1, -2, -3, -4, -13, -14, -15, -16, 9, 10, 11, 12]]
intertwiner: Schur average with the seed matrix unit (1, 2); 2 sqrt(2) U has entries 0, +-1,
    ...+-i (rounding error < 1e-12) with 8 non-zero entries in each row and column;
    ...Gamma17[A].U8 = U8.R_A and U8^dag.U8 = 8 1_16 exactly : True
Schur's lemma: the linear map X -> (Gamma_A X - X R_A)_A has a one-dimensional null space (one
    ...singular value < 1e-10, the next > 0.1) : True
NUM Re(2 sqrt(2) U) : [[0, 1, 0, 0, 0, 0, 1, 0, 0, 1, 0, 0, 0, 0, 1, 0], [0, 0, 0, 1, -1, 0, 0,
    ...0, 0, 0, 0, -1, 1, 0, 0, 0], [0, 0, 1, 0, 0, -1, 0, 0, 0, 0, 1, 0, 0, -1, 0, 0], [1, 0,
    ...0, 0, 0, 0, 0, 1, -1, 0, 0, 0, 0, 0, 0, -1], [0, 0, 0, -1, -1, 0, 0, 0, 0, 0, 0, -1, -1,
    ...0, 0, 0], [0, -1, 0, 0, 0, 0, 1, 0, 0, 1, 0, 0, 0, 0, -1, 0], [-1, 0, 0, 0, 0, 0, 0, 1,
    ...-1, 0, 0, 0, 0, 0, 0, 1], [0, 0, -1, 0, 0, -1, 0, 0, 0, 0, 1, 0, 0, 1, 0, 0], [0, 1, 0,
    ...0, 0, 0, 1, 0, 0, -1, 0, 0, 0, 0, -1, 0], [0, 0, 0, 1, -1, 0, 0, 0, 0, 0, 0, 1, -1, 0,
    ...0, 0], [0, 0, 1, 0, 0, -1, 0, 0, 0, 0, -1, 0, 0, 1, 0, 0], [1, 0, 0, 0, 0, 0, 0, 1, 1,
    ...0, 0, 0, 0, 0, 0, 1], [0, 0, 0, -1, -1, 0, 0, 0, 0, 0, 0, 1, 1, 0, 0, 0], [0, -1, 0, 0,
    ...0, 0, 1, 0, 0, -1, 0, 0, 0, 0, 1, 0], [-1, 0, 0, 0, 0, 0, 0, 1, 1, 0, 0, 0, 0, 0, 0,
    ...-1], [0, 0, -1, 0, 0, -1, 0, 0, 0, 0, -1, 0, 0, -1, 0, 0]]
NUM Im(2 sqrt(2) U) : [[0, 0, 0, -1, -1, 0, 0, 0, 0, 0, 0, 1, 1, 0, 0, 0], [0, 1, 0, 0, 0, 0,
    ...-1, 0, 0, 1, 0, 0, 0, 0, -1, 0], [-1, 0, 0, 0, 0, 0, 0, 1, 1, 0, 0, 0, 0, 0, 0, -1], [0,
    ...0, 1, 0, 0, 1, 0, 0, 0, 0, 1, 0, 0, 1, 0, 0], [0, 1, 0, 0, 0, 0, 1, 0, 0, -1, 0, 0, 0,
    ...0, -1, 0], [0, 0, 0, -1, 1, 0, 0, 0, 0, 0, 0, -1, 1, 0, 0, 0], [0, 0, 1, 0, 0, -1, 0, 0,
    ...0, 0, -1, 0, 0, 1, 0, 0], [-1, 0, 0, 0, 0, 0, 0, -1, -1, 0, 0, 0, 0, 0, 0, -1], [0, 0,
    ...0, -1, -1, 0, 0, 0, 0, 0, 0, -1, -1, 0, 0, 0], [0, 1, 0, 0, 0, 0, -1, 0, 0, -1, 0, 0, 0,
    ...0, 1, 0], [-1, 0, 0, 0, 0, 0, 0, 1, -1, 0, 0, 0, 0, 0, 0, 1], [0, 0, 1, 0, 0, 1, 0, 0,
    ...0, 0, -1, 0, 0, -1, 0, 0], [0, 1, 0, 0, 0, 0, 1, 0, 0, 1, 0, 0, 0, 0, 1, 0], [0, 0, 0,
    ...-1, 1, 0, 0, 0, 0, 0, 0, 1, -1, 0, 0, 0], [0, 0, 1, 0, 0, -1, 0, 0, 0, 0, 1, 0, 0, -1,
    ...0, 0], [-1, 0, 0, 0, 0, 0, 0, -1, 1, 0, 0, 0, 0, 0, 0, 1]]
brane Lorentz group on the brane spinor index, transverse rotations on a: Gamma_mu Gamma_nu = U
    ...(gamma_mu gamma_nu (x) 1) U^dag, Gamma_k Gamma_l = U (1 (x) rho_k rho_l) U^dag, Gamma_mu
    ...Gamma_k = U (gamma_mu gamma5 (x) rho_k) U^dag : True
Spin(4) on the transverse index: rho4567 = -gamma5 (4x4) commutes with rho_k rho_l,
    ...anticommutes with rho_k; (rho_k rho_l + s rho_m rho_n) P_s = 0 and (rho_k rho_l - s
    ...rho_m rho_n) P_s has rank 2 for even (k,l,m,n): 4 = (2,1) + (1,2) : True
chirality: chi = i Gamma_0...Gamma_7 = U (gamma5 (x) rho4567) U^dag; P_L17 = U (P_L (x) P_+ +
    ...P_R (x) P_-) U^dag, each part of rank 4 : True
example psi = U.(u (x) e_1), u = (1, 2, 0, -1): psibar psi = ubar u = 4 and psibar Gamma^mu psi
    ...= ubar gamma^mu u : True
NUM example psi = U.(u (x) e_1), u = (1, 2, 0, -1): psibar psi, psibar Gamma^mu psi (mu = 0..3)
    ...: [4, 6.0, 2.0, 0.0, -4.0]
four brane Dirac spinors: psibar psi' = Sum_a psibar_a psi'_a and psibar Gamma^mu psi' = Sum_a
    ...psibar_a gamma^mu psi'_a (random integer spinors) : True
Kaluza-Klein operator X = gamma5 (x) (rho.q)/b0: X^2 = -|q|^2/b0^2 1_16 and X anticommutes with
    ...gamma_mu (x) 1_4 : True
brane tension: T_mu nu = (2/sqrt|h|) dS/dh^mu nu = f_W h_mu nu for S = -Int sqrt|h| f_W; T^0_0
    ...= f_W, T^1_1 = f_W (p = -f_W, w = -1) : True
smeared sources: tension eps = f_W/V_perp, pa = -eps, pb = 0; brane radiation pa = eps/3, pb =
    ...0; bulk radiation <cos^2> = 3/7, <sin^2> = 4/7: pa = pb = eps/7 : True
normalcy temperature: (T^7/M8^6)/(T^2/M_Pl) = 1 at T* = (M8^6/M_Pl)^(1/5) = 13.3 GeV; V_perp
    ...T^4 T^3/M_Pl^2 = T^7/M8^6 : True
NUM T* = (M8^6/M_Pl)^(1/5) (GeV), f_W/rho_Lambda for f_W = M8^4, 1/(2 pi M8) (GeV^-1) :
    ...(13.264784746562139, 3.9709863911436803e+62, 1.5915494309189534e-05)
the tension is vacuum energy: kappa f_W/V_perp = f_W/M_Pl^2 (V_perp = kappa M_Pl^2); f_W = M8^4
    ...is 3e62 to 5e62 rho_Lambda : True
Green's function of the four transverse directions: (1/y^3) d/dy (y^3 d/dy y^-2) = 0, flux 2
    ...pi^2 y^3 d/dy y^-2 = -4 pi^2; sqrt(kappa f_W) = 1/M8 for f_W = M8^4 : True
boolean checks: 37   failed: 0
VERDICT: ALL CHECKS PASSED
\end{Verbatim}
\end{document}
